

\part{Linear Algebra}
\chapter{Analytic Geometry}
\section{Arithmetic in $n$ dimensions}
We now focus our attention on $n$-dimensional euclidean space, which is modelled by the cartesian product $\bR^n$:
\begin{align*}
	\bR^n = 
	\lr{
	\vx 
	=
	\begin{pmatrix}
		x_1\\
		\vdots\\
		x_n
	\end{pmatrix}
	~
	\mid
	~
	x_1, \hdots, x_n \in \bR}
	=
	\lr{\vx : \lr{1, \hdots, n} \to \bR}
\end{align*}
Elements of $\bR^n$ are known as \tib{vectors}\sidenote{We will see in the next chapter that $\bR^n$ is an example of a \tib{vector space}, including the case $n = 1$ (and $n = 0$, if you want to get a bit pathological). Therefore, formally speaking and only in certain contexts, regular old real (or even rational) numbers can be vectors too. However, when encountering the term "vector" in the wild, it is generally implied that $n \geq 2$.}.
In the context of linear algebra, they are usually written vertically, as \tib{column vectors}.
\newpar
Vectors simultaneously represent both \tib{positions} and \tib{directions} in space - in the same way that the number $5$ is both the point that lies $5$ units away from $0$ on the real number line \tit{and} the distance between the numbers $3$ and $8$ on that same number line, the vector $\binom{1}{2}$ represents both
\begin{itemize}
	\item The point on the cartesian plane that can be reached by going one unit in the direction of one coordinate axis and then two units in the direction of the other axis, \tit{and}
	\item the distance between two points with coordinates $\binom{2}{3}$ and $\binom{3}{5}$, or any other pair of points that are one unit apart in one coordinate direction and two in the other.
\end{itemize}
As is expected for algebraic structures defined on cartesian products, addition is defined \tbf{componentwise}:
\begin{align*}
	\begin{pmatrix}
		x_1\\
		\vdots\\
		x_n
	\end{pmatrix}
	+
	\begin{pmatrix}
		y_1\\
		\vdots\\
		y_n
	\end{pmatrix}
	=
	\begin{pmatrix}
		x_1 + y_1\\
		\vdots\\
		x_n + y_n
	\end{pmatrix}
\end{align*}
\newpar
Given a real number $c$, we can scale elements of $\bR^n$ by that number:
\begin{align*}
	c \cdot \vx
	=
	c \cdot 
	\begin{pmatrix}
		x_1\\
		\vdots\\
		x_n
	\end{pmatrix}
	=
	\begin{pmatrix}
		c \cdot x_1\\
		\vdots\\
		c \cdot x_n
	\end{pmatrix}
\end{align*}
The factor $c$ is known as a \tib{scalar}, and the operator $\cdot : \bR \times \bR^n \to \bR^n$ as \tib{scalar multiplication}.
For the sake of readability, we differentiate scalars from vectors by writing a little arrow on top of a vector.
\newpar 
Arguably, it is these two operations - vector addition and scalar multiplication - that define linear algebra - next chapter, we will abstractly define \tib{vector spaces} to be those structures where addition and scalar multiplication are well-behaved and the scalars come from a field, and \tib{modules} to be their cousins whose scalars come from a ring.
\newpar
In general, there is no well-behaved\sidenote{In this case, our criterion for multiplication being well-behaved is that multiplicative inverses should exist, i.e. that division should be possible.} vector multiplication that works for any $n$. Later on, we will prove this observation in the form of \tib{Frobenius' Theorem}, which roughly states:
\begin{importanttheorem}{Frobenius' Theorem}{}
	The only dimensions in which an associative multiplication exists on $\bR^n$ are $\bR^1$, i.e. $\bR$ itself, $\bR^2$, where multiplication is given by treating two-dimensional vectors as complex numbers, and $\bR^4$, where multiplication is given treating four-dimensional vectors as \tib{quaternions}.
\end{importanttheorem} 
We will take a closer look at the complex numbers $\bC \cong \bR^2$ and the quaternions $\bH \cong \bR^4$ in the following sections. In particular, we will see that quaternion multiplication isn't commutative, so if we require multiplication to be commutative, our options shrink even further, leaving just $\bR$ and $\bC \cong \bR^2$.

\clearpage

\section{The Inner Product: Angles and Lengths in $\bR^n$}
Another operation that is of central importance to euclidean geometry is the \tib{inner product}, also known as the \tib{dot product} or \tib{scalar product}.
\begin{importantdefinition}{Euclidean Inner Product}{}
	Given two vectors $\vx, \vy \in \bR^n$, their \tib{Euclidean inner product}, written $\scalar{\vx}{\vy}$ or simply $\vx \cdot \vy$, is the sum of the products of the components of $\vx$ with the components of $\vy$:
	\begin{align*}
		\scalar{\vx}{\vy} = \sum_{i = 1}^n x_i \cdot y_i
	\end{align*}
\end{importantdefinition}
\begin{importanttheorem}{Properties of Inner Products}{}
	The Euclidean inner product has the following properties:
	\begin{enumerate}
		\item \tib{(Positive definiteness):} For every $\vx \in \bR^n$, we have $\scalar{\vx}{\vx} = 0$ if and only if $\vx = \vzero$.
		\item \tib{(Linearity in the first argument):} For every $\vx, \vy, \vz \in \bR^n$ and $a,b \in \bR$, we have:
		\begin{align*}
			\scalar{ax + by}{z} = a \scalar{x}{z} + b \scalar{y}{z}
		\end{align*}
		\item \tib{(Symmetry):} For every $\vx, \vy \in \bR^n$, we have:
		\begin{align*}
			\scalar{x}{y} = \scalar{y}{x}
		\end{align*}
	\end{enumerate}
\end{importanttheorem}
\begin{exercise}
	Show that the Euclidean inner product fulfills the following version of the binomial formula:
	\begin{align*}
		\scalar{\vx + \vy}{\vx + \vy} = \scalar{\vx}{\vx} + 2 \scalar{\vx}{\vy} + \scalar{\vy}{\vy}\\
		\scalar{\vx - \vy}{\vx - \vy} = \scalar{\vx}{\vx} - 2 \scalar{\vx}{\vy} + \scalar{\vy}{\vy}
	\end{align*}
\end{exercise}
The importance of inner products comes from the fact they let us define \tib{lengths} and \tib{angles}.
\begin{importantdefinition}{Euclidean Norm}{}
	The length of a vector $\vx \in \bR^n$, also known as its \tib{Euclidean norm}, written as $\norm{\vx}_2$, $\norm{\vx}$, or just $\abs{x}$, is given by:
	\begin{align*}
		\norm{\vx}_2 = \sqrt{\scalar{\vx}{\vx}} = \sqrt{x_1^2 + \hdots + x_n^2}
	\end{align*}
\end{importantdefinition}
Note that for $n = 1$, we simply get $\norm{x}_2 = \sqrt{x^2} = \abs{x}$, so any result concerning the Euclidean norm on $\bR^n$ also holds for the absolute value function.
\begin{importanttheorem}{Some properties of Norms}{}
	The Euclidean norm has the following properties:
	\begin{enumerate}
		\item \tib{(Non-negativity):} For every $\vx \in \bR^n$, $\norm{\vx}_2 \geq 0$.
		\item \tib{(Positive definiteness):} For every $\vx \in \bR^n$, $\norm{\vx}_2 = 0$ if and only if $\vx = \vzero$.
		\item \tib{(Absolute homogeneity):} For every $c \in \bR$ and $\vx \in \bR^n$, we have
		\begin{align*}
			\norm{c \cdot \vx}_2 = \abs{c} \cdot \norm{\vx}_2
		\end{align*}
	\end{enumerate}
\end{importanttheorem}
The little subscript $2$ in $\norm{\vx}_2$ differentiates the Euclidean norm from the many other norms we will meet over time - in general, for any $p \in \bN_1$, $\norm{\vx}_p$ denotes the function
\begin{align*}
	\norm{\vx}_p = \sqrt[p]{\abs{x_1}^p + \hdots + \abs{x_n}^p} \qquad\lr(= \sqrt[p]{\sum_{i = 1}^n \abs{x_i}^p}),
\end{align*}
known as the \tib{$p$-norm}. For now, we only really care about the Euclidean norm, but other values of $p$ will start coming in handy once we start applying tools from analysis to vector spaces. The limiting case $p \to \infty$ is particularly useful: We will find that $$\norm{\vx}_\infty = \max \abs{x_i}.$$
\begin{importanttheorem}{Cauchy-Schwartz Inequality}{cauchyschwartzspecific}
	Let $\vx, \vy \in \bR^n$. Then we have
	\begin{align*}
		\abs{\scalar{\vx}{\vy}} \leq \norm{\vx}_2 \cdot \norm{\vy}_2
	\end{align*}
\end{importanttheorem}
\begin{proof}
	\begin{enumerate}
		\item Assume $\norm{\vx}_2 = 0$. Then we immediately get 
		\begin{align*}
			\scalar{\vx}{\vy} = 0 = \norm{\vx}_2 \cdot \norm{\vy}_2.
		\end{align*}
		\item Assume $\norm{\vx}_2 = \norm{\vy}_2 = 1$. Then we have $\norm{\vx}_2^2 + \norm{\vy}_2^2 = 2$, and thus:
		\begin{align*}
			1 + \scalar{\vx}{\vy}
			&= \frac{1}{2} (\norm{\vx}_2^2 + \norm{\vy}_2^2) + \scalar{\vx}{\vy}\\
			&= \frac{1}{2} (\norm{\vx}_2^2 + 2\scalar{\vx}{\vy} + \norm{\vy}_2^2)\\
			&= \frac{1}{2} \norm{\vx + \vy}_2^2\\
			&\geq 0
		\end{align*}
		Thus we get $\scalar{\vx}{\vy} \geq -1$. Applying the same reasoning to $1 - \scalar{\vx}{\vy}$, we get 
		\begin{align*}
			1 - \scalar{\vx}{\vy} = \frac{1}{2} \norm{\vx - \vy}_2^2 \geq 0,
		\end{align*}
		i.e. $\scalar{\vx}{\vy} \leq 1$. We thus have 
		\begin{align*}
			\abs{\scalar{\vx}{\vy}} \leq 1 = \norm{\vx}_2 \cdot \norm{\vy}_2,
 		\end{align*}
	 	as desired.
	 	\item Now, let $\vx, \vy \in \bR^n \setminus \lr{\vzero}$ be arbitrary. Then $\frac{\vx}{\norm{\vx}_2}$ and $\frac{\vy}{\norm{\vy}_2}$ both have norm $1$, letting us simply use our previous result:
	 	\begin{align*}
	 		\abs{\scalar{\vx}{\vy}}
	 		&= \abs{\norm{\vx} \cdot \norm{\vy} \cdot \scalar{\frac{\vx}{\norm{\vx}}}{\frac{\vy}{\norm{\vy}}}}\\
	 		&\leq \abs{\norm{\vx} \cdot \norm{\vy} \cdot \norm{\frac{\vx}{\norm{\vx}}} \cdot \norm{\frac{\vy}{\norm{\vy}}}}\\
	 		&= \norm{\vx} \cdot \norm{\vy}
	 	\end{align*}
	\end{enumerate}
\end{proof}
\begin{importantcorollary}{Triangle Inequality}{}
	For every $\vx, \vy \in \bR^n$, we have
	\begin{align*}
		\norm{\vx + \vy}_2 \leq \norm{\vx}_2 + \norm{\vy}_2
	\end{align*}
\end{importantcorollary}
We had previously formalized the concept of an 'angle between two lines' via the complex exponential function, and then defined the trigonometric functions $\sin x$ and $\cos x$ as the complex and real parts of $e^{ix}$, respectively. If we take the cosine function as a given, then scalar products provide a convenient alternate way of defining angles, which works in any dimension:
\begin{importantdefinition}{Inner angle between two vectors}{}
	Given two vectors $\vx,\vy \in \bR^n \setminus \lr{\vzero}$, we define the \tib{inner angle} between them to be:
	\begin{align*}
		\angle{\vx}{\vy} = \arccos \lr( \frac{\scalar{\vx}{\vy}}{\norm{\vx}_2 \cdot \norm{\vy}_2})
	\end{align*}
\end{importantdefinition}
Note that Cauchy-Schwartz gives us $\frac{\scalar{\vx}{\vy}}{\norm{\vx}_2 \cdot \norm{\vy}_2} \in [-1,1]$, showing that this expression actually is well-defined.
\begin{exercise}
	Verify that we have:
	\begin{enumerate}
		\item $\angle{\vx}{\vy} = 0$ if and only if there exists $c \in \bR_{> 0}$ such that $\vx = c \vy$.
		\item $\angle{\vx}{\vy} = \frac{\pi}{2}$ if and only if $\scalar{x}{y} = 0$.
		\item $\angle{\vx}{\vy} = \pi$ if and only if there exists $c \in \bR_{> 0}$ such that $\vx = -c \vy$.
	\end{enumerate}
\end{exercise}
Notably, this suggests that our definition is sensible even in the case $n = 1$, since any pair of non-zero vectors fulfills either condition 1 or 2, implying that the angle between two vectors on the real number line is either $0$ or $\pi$.
\newpar
These inner angles match our earlier definition of angles in terms of complex numbers, as long as we are content with not being able to meaningfully talk about negative angles or angles greater than $\pi$ anymore:
\begin{exercise}
	Earlier in this book, we defined angles in terms of complex numbers. Given a complex number $y = r e^{i \phi}$ we saw that the corresponding real vector is given by $\vy = \binom{r \cos \phi}{r \sin \phi}$. 
	\newpar
	Show that our two definitions match as long as $\phi \in [0, \pi]$, i.e. that given a vector $\vx = \binom{x}{0}$ lying on the $x$-axis, we have
	\begin{align*}
		\angle{\vx}{\vy} = \phi
	\end{align*}
	(Hint: Remember $\cos^2 \phi + \sin^2 \phi = 1$)
\end{exercise}
We have to be content with this restricted interval $[0,\pi]$, since the cosine function isn't injective on larger intervals, meaning the arccosine would no longer be well-defined.
\newpar
More generally, any definition of an 'angle between two arbitrary vectors' will suffer from this - if we only specify two vectors, it's left ambiguous which side we should measure the angle on.
\newpar 
This problem can be solved by finding a way to take a tuple $(\vx_1, \hdots, \vx_n)$ of vectors and recognize whether it is 'positively ordered' or 'negatively ordered' - in $\bR^2$, this would be asking whether the second vector lies 'clockwise' or 'counterclockwise' from the first, while in $\bR^3$ the question is whether the vectors are 'left-handed' or 'right-handed'. A map that gives us this information is known as an \tib{orientation}. We will formalize this idea later, when we have introduced the necessary tools, those being vector spaces, their \tib{bases}, and \tib{determinant functions}.
\newpar
In the complex case, we could dodge this issue because one of our vectors, call it $\vx$, was always part of the $x$-axis, and thus we could obtain a canonical orientation by declaring that the angle between $\vx$ and $\vy$ was positive if $y_2 > 0$ and negative if $y_2 < 0$.
\begin{exercise}
	Prove the following theorems of classical Euclidean geometry:
	\begin{enumerate}
		\item The \tib{Law of Cosines}:
		\begin{align*}
			\norm{\vx-\vy}_2^2 = \norm{\vx}_2^2 + \norm{\vy}^2 - 2 \norm{\vx}_2 \norm{\vy}_2 \cdot \cos(\angle{\vx}{\vy})
		\end{align*}
		\item The \tib{Pythagorean Theorem}:
		\begin{align*}
			\angle{\vx}{\vy} = \frac{\pi}{2} \implies \norm{\vx-\vy}_2^2 = \norm{\vx}_2^2 + \norm{\vy}^2
		\end{align*}
	\end{enumerate}
\end{exercise}
\iffalse
In the-dimensional case $x \in \bR$, $y = r \cdot e^{\phi i} \in \bC$, corresponding to the real vectors $\vx = \binom{x}{0}$, $\vy = \binom{r \cos \phi}{r \sin \phi}$, this definition matches our previous complex definition as long as we restrict $\phi$ to the interval $[0,\pi]$. In this case, we have:
\begin{align*}
	\angle (\vx, \vy) 
	&= \arccos \lr(\frac{\scalar{\vx}{\vy}}{\norm{\vx}_2 \cdot \norm{\vy}_2})\\
	&= \arccos \lr(\frac{(x \cdot r \cos \phi) + (0 \cdot r \sin \phi)}{\sqrt{x^2 + 0^2} \cdot \sqrt{(r \cos(\phi))^2 + (r\sin(\phi))^2}})\\
	&= \arccos \lr(\frac{xr\cos \phi}{xr})\\
	&= \arccos(\cos \phi)\\
	&= \phi
\end{align*}
Where the third step made use of the trigonometric Pythagorean theorem $\cos(\phi)^2 + \sin(\phi)^2 = 1$. 
\fi

\clearpage

\section{The Cross Product}
The \tib{cross product} is an operation $\times : \bR^3 \times \bR^3 \to \bR^3$, such that:
\begin{enumerate}
	\item $\vu \times \vv$ is orthogonal to both $\vu$ and $\vv$,
	\item the length of $\vu \times \vv$ is the area of the parallelogram spanned by $\vu$ and $\vv$:
	\begin{align*}
		\norm{\vu \times \vv}_2 = \norm{\vu}_2 \norm{\vv}_2 \sin(\angle{\vu}{\vv})
	\end{align*}
\end{enumerate}

\newpar
We also declare that our cross product should be \tib{anticommutative:} given two vectors, we want to be able to change the sign of the cross product by switching the order in which we multiply:
\begin{align*}
	\vv \times \vw := -\vw \times \vv
\end{align*}
Notice how this immediately implies
\begin{align*}
	\vv \times \vv &= - (\vv \times \vv)\\
	\implies \vv \times \vv &= 0
\end{align*}
Finally, notice that our geometric intuition immediately tells us that our product should be \tib{linear} and \tib{distributive over vector addition}. 
\newpar
We can now derive an explicit formula for this operation on our own. 
First, note that the standard basis
\begin{align*}
	\ve_1 = \trinom{1}{0}{0}, \ve_2 = \trinom{0}{1}{0}, \ve_3 = \trinom{0}{0}{1},
\end{align*}
consists of orthogonal vectors of length one, so we can start by declaring
\begin{align*}
	\ve_1 \times \ve_2 &:= \ve_3\\
	\ve_2 \times \ve_3 &:= \ve_1\\
	\ve_3 \times \ve_1 &:= \ve_2
\end{align*}
Note that once again, the direction in which these vectors point depends on a somewhat arbitrary choice of orientation - we could invert all of the directions and still end up with a valid definition.
\newpar
These observations are already enough to calculate the cross product of two arbitrary vectors: Using distributivity, we get:
\begin{align*}
	\trinom{v_1}{v_2}{v_3} \times \trinom{w_1}{w_2}{w_3}
	&=
	(v_1 \ve_1 + v_2\ve_2 + v_3\ve_3) \times (w_1 \ve_1 + w_2\ve_2 + w_3\ve_3)\\
	&=
	v_1w_1 (\ve_1 \times \ve_1)\\
	&+
	v_1w_2 (\ve_1 \times \ve_2)\\
	&+
	v_1w_3 (\ve_1 \times \ve_3)\\\\
	&+
	v_2w_1 (\ve_2 \times \ve_1)\\
	&+
	v_2w_2 (\ve_2 \times \ve_2)\\
	&+
	v_2w_3 (\ve_2 \times \ve_3)\\\\
	&+
	v_3w_1 (\ve_3 \times \ve_1)\\
	&+
	v_3w_2 (\ve_3 \times \ve_2)\\
	&+
	v_3w_3 (\ve_3 \times \ve_3)
\end{align*}
To which we can apply anticommutativity and our formulas for multiplying the basis vectors to get:
\begin{align*}
	&\trinom{v_1}{v_2}{v_3} \times \trinom{w_1}{w_2}{w_3}\\\\
	&=
	v_1w_1 \cdot \vzero\\
	&+
	v_1w_2 \cdot \ve_3\\
	&+
	v_1w_3 \cdot (- \ve_2)\\\\
	&+
	v_2w_1 \cdot (- \ve_3)\\
	&+
	v_2w_2 \cdot \vzero\\
	&+
	v_2w_3 \cdot \ve_1\\\\
	&+
	v_3w_1 \cdot \ve_2\\
	&+
	v_3w_2 \cdot (-\ve_1)\\
	&+
	v_3w_3 \cdot \vzero\\\\
	&= \trinom{v_2w_3 - v_3w_2}{v_3w_1 - w_1v_3}{v_1w_2 - v_2w_1}
\end{align*}
\begin{importantcorollary}{Cross Product}{}
	The cross product $\vv \times \vw$ of two vectors $\vv, \vw \in \bR^3$ returns a vector whose direction is orthogonal to both $\vv$ and $\vw$ and whose length equals the area of the parallelogral spanned between them. Given coordinate representations of $\vv$ and $\vw$, the cross product can be calculated via the following formula:
	\begin{align*}
		\trinom{v_1}{v_2}{v_3} \times \trinom{w_1}{w_2}{w_3}
		=
		\trinom{v_2w_3 - v_3w_2}{v_3w_1 - w_1v_3}{v_1w_2 - v_2w_1}
	\end{align*}
\end{importantcorollary}
\newpar
\begin{proposition}
	We have
	\begin{align*}
		\scalar{\vu \times \vv}{\vw \times \vt} = \scalar{\vu}{\vw}\scalar{\vv}{\vt} - \scalar{\vu}{\vt}\scalar{\vv}{\vw}.
	\end{align*}
\end{proposition}
\begin{proposition}
	We have
	\begin{align*}
		\scalar{\vu \times \vv}{\vw} = \scalar{\vv \times \vw}{\vu} = \scalar{\vw \times \vu}{\vv}
	\end{align*}
	This expression is known as the \tib{box product} of $\vu, \vv, \vw$. Its absolute value corresponds to the volume of the parallelepiped spanned by the three vectors.
\end{proposition}
\begin{proposition}
	We have 
	\begin{align*}
		(\vu \times \vv) \times \vw = \scalar{u}{w} \cdot v - \scalar{v}{w} \cdot u = w \times (v \times u).
	\end{align*}
	This result is known as \tib{Grassmann's Identity}.
\end{proposition}
\begin{proposition}
	We have
	\begin{align*}
		(\vu \times \vv) \times \vw + (\vv \times \vw) \times \vu + (\vw \times \vu) \times \vv = 0
	\end{align*}
	This result is known as \tib{Jacobi's Identity.}
\end{proposition}
\clearpage

\section{The Quaternions $\bH$}
\begin{importantdefinition}{Quaternions}{}
	The Quaternions $\bH$ consist of the set $\bR^4$, with the following special notation:
	\begin{align*}
		\vec 1 &:= (1,0,0,0)\\
		i &:= (0,1,0,0)\\
		j &:= (0,0,1,0)\\
		k &:= (0,0,0,1)
	\end{align*}
	Such that in addition to scalar multiplication and vector addition, which work as usual, we have an additional vector multiplication operation $\cdot : \bR^4 \times \bR^4 \to \bR^4$ that behaves as follows:
	\begin{enumerate}
		\item $\vec 1$ is a neutral element of vector multiplication: 
		\begin{align*}
			\forall x &\in \bH:\\ 
			\vec 1 \cdot x &= x \cdot \vec 1 = x
		\end{align*}
		\item Vector multiplication is associative:
		\begin{align*}
			\vx(\vy\vz) = (\vx\vy)\vz
		\end{align*}
		\item Multiplicative inverses exist:
		\begin{align*}
			\forall \vx &\in \bH_{\neq 0}:\\
			\exists \vx^{-1} &\in \bH:\\
			\vx \cdot \vx^{-1} &= \vx^{-1} \cdot \vx = \vec 1
		\end{align*}
		\item Vector multiplication distributes over vector addition, from both sides:
		\begin{align*}
			\forall \vx, \vy, \vz &\in \bH:\\ 
			\vx \cdot (\vy + \vz) &= \vx  \vy + \vx \vz\\
			(\vx + \vy) \cdot \vz &= \vx \vz + \vy \vz´
		\end{align*}
		\item Vector multiplication is compatible with scalar multiplication:
		\begin{align*}
			\forall a,b &\in \bR,
					\vx,\vy \in \bH:\\
			a\vx \cdot b\vy &= (ab)(\vx\vy)
		\end{align*}
		\item The unit vectors fulfill the \tib{Quaternion multiplication formula}:
		\begin{align*}
			i^2 = j^2 = k^2 = ijk = -\vec 1
		\end{align*}
	\end{enumerate}
\end{importantdefinition}
Apart from the Quaternion multiplication formula, which is the defining feature of the Quaternions, all of our other demands together state that $\bH$ should be an \tib{associative real division algebra over $\bR^4$}.
\newpar
One very important consequence of this definition is that vectors of the form $\vec a := (a,0,0,0)$ behave exactly like real numbers - let $\vec x \in \bH$, then we have:
\begin{align*}
	\vec a \vec x = (a \cdot \vec1) \cdot \vx = a \cdot (\vec 1 \cdot \vx) = a \vx
\end{align*}
For this reason, it is common to write $\vec 1$ simply as $1$, and to then proceed as we did for the complex numbers and write:
\begin{align*}
	(a,b,c,d) = a + bi + cj + dk
\end{align*}
We can now already get most of the way to deriving multiplication formula for quaternions, by simply applying distributivity:
\begin{align*}
	(v_1 + v_2i + v_3j + v_4k) &\cdot (u_1 + u_2i + u_3j + u_4k)\\
	&= v_1u_1 + (v_1u_2)i + (v_1u_3)j + (v_1u_4)k\\
	&+ (v_2u_1)i + (v_2u_2)i^2 + (v_2u_3)ij + (v_2u_4)ik\\
	&+ (v_3u_1)j + (v_3u_2)ji + (v_3u_3)j^2 + (v_3u_4)jk\\
	&+ (v_4u_1)k + (v_4u_2)ki + (v_4u_3)kj + (v_4u_4)k^2\\
\end{align*}
Now, using $i^2 + j^2 + k^2 = ijk = -1$, we can simplify slightly:
\begin{align*}
	(v_1 + v_2i + v_3j + v_4k) &\cdot (u_1 + u_2i + u_3j + u_4k)\\
	&= 
		v_1u_1 + (v_1u_2)i + (v_1u_3)j + (v_1u_4)k\\
	&+ (v_2u_1)i + (v_2u_2)i^2 + (v_2u_3)ij + (v_2u_4)ik\\
	&+ (v_3u_1)j + (v_3u_2)ji + (v_3u_3)j^2 + (v_3u_4)jk\\
	&+ (v_4u_1)k + (v_4u_2)ki + (v_4u_3)kj + (v_4u_4)k^2\\
	&= 
		v_1u_1 + (v_1u_2)i + (v_1u_3)j + (v_1u_4)k\\
	&+ (v_2u_1)i - (v_2u_2) + (v_2u_3)ij + (v_2u_4)ik\\
	&+ (v_3u_1)j + (v_3u_2)ji - (v_3u_3) + (v_3u_4)jk\\
	&+ (v_4u_1)k + (v_4u_2)ki + (v_4u_3)kj - (v_4u_4)\\
\end{align*}
\newpar
All that's left is to figure out what the products of $i,j$ and $k$ should be.
\newpar
We can derive $ij$ and $ji$ as follows:
\begin{enumerate}
	\item Start with $ijk = -1$.
	\item Multiply $k$ from the right to get $ijk^2 = -k$, i.e. $-ij = -k$, i.e. $ij = k$.
	\item Multiply $ij$ from the left to get $ijij = ijk = -1$.
	\item Multiply $i$ from the left and $j$ from the right to get $i(ijij)j = -ij$.
	\item Simplify the left side to get $i(ijij)j = (ii)(ji)(jj) = ji$, i.e:
	\begin{align*}
		ij = -ji = k
	\end{align*}
\end{enumerate}
\begin{exercise}
	Show that we also have $jk = -kj = i$ and $ki = -ik = j$.
\end{exercise}
This means that Quaternion multiplication isn't commutative!
\newpar
Our final multiplication table for $1,i,j$ and $k$ is thus:
\begin{align*}
	\begin{array}{c|cccc}
		\cdot & 1 & i & j & k\\\hline
		1 & 1 & i & j & k\\
		i & i &-1 & k &-j\\
		j & j &-k &-1 & i\\
		k & k & j &-i &-1\\
	\end{array}
\end{align*}
\begin{exercise}
	Verify that this multiplication table doesn't violate associativity.
\end{exercise}
Entering these new results into our partial multiplication formula, we get:
\begin{align*}
	(v_1 + v_2i + v_3j + v_4k) &\cdot (u_1 + u_2i + u_3j + u_4k)\\
	&= 
	v_1u_1 + (v_1u_2)i + (v_1u_3)j + (v_1u_4)k\\
	&+ (v_2u_1)i - (v_2u_2) + (v_2u_3)ij + (v_2u_4)ik\\
	&+ (v_3u_1)j + (v_3u_2)ji - (v_3u_3) + (v_3u_4)jk\\
	&+ (v_4u_1)k + (v_4u_2)ki + (v_4u_3)kj - (v_4u_4)\\
	&=
	v_1u_1 + (v_1u_2)i + (v_1u_3)j + (v_1u_4)k\\
	&+ (v_2u_1)i - (v_2u_2) + (v_2u_3)k - (v_2u_4)j\\
	&+ (v_3u_1)j - (v_3u_2)k - (v_3u_3) + (v_3u_4)i\\
	&+ (v_4u_1)k + (v_4u_2)j - (v_4u_3)i - (v_4u_4)\\
\end{align*}
which we can finally simplify to get:
\begin{align*}
	(v_1 + v_2i + v_3j + v_4k) \cdot (u_1 + u_2i + u_3j + u_4k)\\
	=~ (v_1u_1 - v_2u_2 - v_3u_3 - v_4u_4)\\
	+~ (v_1u_2 + v_2u_1 + v_3u_4 - v_4u_3)i\\
	+~ (v_1u_3 - v_2u_4 + v_3u_1 + v_4u_2)j\\
	+~ (v_1u_4 + v_2u_3 - v_3u_2 + v_4u_1)k
\end{align*}
Now, this is obviously a long and unreasonably complicated formula - is there any way to think of it differently? The complex numbers had a very nice and intuitive geometric definition, maybe the quaternions do too?
\newpar
Well, the key trick here is to split the underlying set $\bR^4$ into $\bR \times \bR^3$, thinking of each Quaternion as consisting of a \tib{real part} $a \in \bR$ and an \tib{imaginary part} $\vv \in \bR^3$. Our multiplication formula now reads:
\begin{align*}
	(a, \vv) \cdot (b, \vu)\\
	=~ (ab - v_1u_1 - v_2u_2 - v_3u_3)\\
	+~ (au_1 + bv_1 + v_2u_3 - v_3u_2)i\\
	+~ (au_2 - v_1u_3 + bv_2 + v_3u_1)j\\
	+~ (au_3 + v_1u_2 - v_2u_1 + bv_3)k
\end{align*}
which you might notice consists exclusively of familiar vector operations:
\begin{enumerate}
	\item $ab - v_1u_1 - v_2u_2 - v_3u_3 = ab - \scalar{\vv}{\vu}$
	\item $\trinom{au_1}{au_2}{au_3} = a\vu$, $\trinom{bv_1}{bv_2}{bv_3} = b\vv$
	\item $\trinom{v_2u_3-v_3u_2}{v_1u_3-v_3u_1}{v_1u_2-v_2u_1} = \vv \times \vu$
\end{enumerate}
Putting everything together, we can now write out our multiplication formula in a much nicer form:
\begin{align*}
	(a, \vv) \cdot (b, \vu) =
	(ab - \scalar{\vv}{\vu}, a\vu + b\vv + \vv \times \vu),
\end{align*} 
\begin{exercise}
	Verify that quaternion multiplication is associative and distributes over vector addition.
\end{exercise}
\begin{exercise}
	Let $p = (a,\vv) \in \bH$. Then $pq = qp$ holds for all $q \in \bH$ if and only if $\vv \neq 0$.
\end{exercise}
Similar to complex numbers, we can define \tib{quaternion conjugation}:
\begin{importantdefinition}{Quaternion Conjugation}{}
	Just like for $\bC$, we define the \tib{conjugate} of a quaternion $(a, \vv)$ to be the quaternion
	\begin{align*}
		\ol{(a, \vv)} = (a, -\vv)
	\end{align*}
\end{importantdefinition}
An immediate application of quaternion conjugation is finding a formula for the inverse of a quaternion, which ones again ends up working out very similarly to the complex case -
We have
\begin{align*}
	\ol{(a, \vv)} \cdot (a, \vv) 
	&= (a, -\vv) \cdot (a, \vv)\\
	&= (a^2 - \scalar{-\vv}{\vv}, a\vv - a\vv + (-\vv \times \vv))\\
	&= \lr(a^2 + \norm{\vv}^2, \vzero) \in \bR
\end{align*}
And thus we can simply divide by this value to get our inverse formula:
\begin{importantcorollary}{Inverse Quaternion}{}
	Let $(a, \vv) \in \bH_{\neq 0}$. Then $(a, \vv)$ has a multiplicative inverse, which is given by:
	\begin{align*}
		(a, \vv)^{-1} = \frac{(a,-\vv)}{a^2 + \norm{v}^2}.
	\end{align*}
\end{importantcorollary}
\begin{definition}
	Let $q = (a, \vv) \in \bH$. Then
	\begin{align*}
		\abs{q} := \sqrt{\ol q \cdot q} = a^2 + v_1^2 + v_2^2 + v_3^2 = \norm{q}_2
	\end{align*}
\end{definition}
Note how much this once again resembles both the real case $\abs{x} = \sqrt{x^2} = \norm{x}_2$ and the complex case $\abs{z} = \sqrt{\ol z z} = \norm{z}_2$.
\clearpage
\section{Quaternions and $\bR^3$}
\begin{importanttheorem}{Quaternions encode rotations}{}
	Let $\norm{\vv} = 1$ and
	\begin{align*}
		q = (\cos \phi, \vv \sin \phi) \in \bH
	\end{align*} 
	Then $q w \ol q \in \bR^3$ is the vector obtained by rotating $w$ around the axis $\lr{\lambda\vv : \lambda \in \bR}$ by the angle $2\phi$.
\end{importanttheorem}
Since $2\pi$ corresponds to a full rotation, $\phi$ and $\phi + \pi$ describe the same rotation. Therefore, there are two different quaternions $q$ and $-q$ associated with each rotation - in topological terms, we say that the quaternions form a \tib{double cover} of the space of all rotations. 
\newpar
This phenomenon is actually closely related to the theory of \tib{spin} in quantum physics. All known quantum particles that constitute ordinary matter, including protons, neutrons, electrons, neutrinos, and quarks, are \tib{spin-$\frac{1}{2}$}, meaning that they exhibit a curious behavior where a rotation of $360$° leaves them in a different state than they were before, from which they return to their original state only after being rotated by another $360$° around the same axis. This lets us mathematically model particles as quaternions, such that a $360$° rotation transforms a particle from a state $q \in \bH$ to the "rotationally equivalent" state $-q \in \bH$.
\begin{corollary}
	For every isometry $F : \bR^3 \to \bR^3$, there exist $u \in \bR^3$ and $q \in \bH$ such that we have either
	\begin{enumerate}
		\item $F(w) = u + q w \ol q$, if $F$ preserves orientation, or
		\item $F(w) = u + q \ol w \ol q$, if $F$ doesn't preserve orientation.
	\end{enumerate}
\end{corollary}
This representation turns out to be a lot more computationally efficient than other representations of the same rotation. Other, more intuitively "trivial" representations of rotation also run into issues such as \tib{gimbal lock}, where they may get "stuck" in certain orientations.
\newpar
This makes quaternions surprisingly practical in many applied contexts, including game development and animation\footnote{For an introduction to quaternions in game development, I highly recommend \href{https://www.youtube.com/watch?v=PMvIWws8WEo}{this excellent talk by Freya Holmer.}}, robotics, and three-dimensional image processing.
\chapter{Linear Maps}
\section{Modules and Vector Spaces}
\begin{importantdefinition}{Modules}{}
	Let $(R, +, \cdot)$ be a Ring. Then a (left) \tib{$R$-Module} consists of an abelian Group $(M, +)$ equipped with a \tib{scalar multiplication} $. : R \times M \to M$, i.e. scalars are elements of $R$ and "vectors" are elements of $M$, such that:
	\begin{enumerate}
		\item Ring multiplication associates with scalar multiplication:
		\begin{align*}
			(r \cdot s).m = r . (s . m)
		\end{align*}
		\item Scalar multiplication distributes over module addition:
		\begin{align*}
			r . (m + n) = r.m + r.n
		\end{align*}
		\item Scalar multiplication distributes over scalar addition:
		\begin{align*}
			(r + s).m = r.m + s.m 
		\end{align*}
	\end{enumerate}
	If $R$ is unital, then we say that $M$ is unital if the identity of ring multiplication is the identity of scalar multiplication:
		\begin{align*}
			1.m = m
		\end{align*}
	
\end{importantdefinition}
A \tib{right $R$-module} is defined similarly, but with a scalar multiplication $. : M \times R \to M$, which is applied from the right instead of from the left. 
\begin{exercise}
	Show that the distinction between right $R$-modules and left $R$-modules is only relevant if $R$ is noncommutative, i.e. that any left $R$-module $M$ over a commutative ring can be turned into a right $R$-module by defining $m.r = r.m$ and vice versa.
\end{exercise}
To keep this text in line with familiar conventions of vector multiplication in $\bR^n$, whenever we talk about a module without explicitly specifiying whether scalar multiplication is applied from the right or from the left, it is implied that it is applied from the left.
\pagebreak
\begin{exercise}
	Let $M$ be a module over $R$. Let $m \in M$ and $r \in R$. show that:
	\begin{enumerate}
		\item $0_R . m = 0_M$
		\item $r . 0_M = 0_M$
		\item $(-r).m = r.(-m) = -(r . m)$
	\end{enumerate}
\end{exercise}
\begin{exercise}
	Verify that:
	\begin{enumerate}
		\item Every ring $R$ forms a module over itself, where scalar multiplication is given by ring multiplication.
		\item $\bR^n$ forms a module over $\bR$, with scalar multiplication defined the usual way. More generally, given any ring $R$, the cartesian product $R^n$ forms an $R$-module where scalar multiplication is given by ring multiplication applied componentwise.
		\item Any abelian group $A$ can be turned into a non-unital module over any ring $R$ by equipping it with the trivial scalar multiplication
		\begin{align*}
			r.a = 0_A
		\end{align*}
	\end{enumerate}
\end{exercise}
\begin{importantdefinition}{Vector Spaces}{}
	A module over a field (or skew field) $\Bbbk$ is known as a \tib{$\Bbbk$-vector space}.
\end{importantdefinition}
\begin{importantdefinition}{Linear Maps}{}
	A \tib{module homomorphism} is a structure-preserving map $F : M \to N$ between two modules over the same ring $R$, i.e. given $m_1,m_2 \in M$ and $r \in R$, we have that:
	\begin{enumerate}
		\item $F$ is \tib{additive}, i.e. it is compatible with module addition: 
		$$F(m_1 +_M m_2) = F(m_1) +_N F(m_2)$$
		\item $F$ is \tib{homogenous}, i.e. it is compatible with scalar multiplication: $$F(r._M m) = r ._N F(m)$$
	\end{enumerate}
	Module homomorphisms are also known as \tib{linear maps}.
\end{importantdefinition}
As the name suggests, the main goal of linear algebra is to generalize the concept of a \tib{linear map}, and to figure out properties and applications of these maps.
\begin{exercise}
	Show that any linear map $F : M \to N$ fulfills
	\begin{align*}
		F(0_M) = 0_N
	\end{align*}
\end{exercise}
\begin{exercise}
	Show that any linear map $F : R \to R$ from a unital ring into itself can be represented as multiplication by a constant, i.e. $F(r) = a.r = a \cdot_R r$ for some $a \in R$.
\end{exercise}
\begin{exercise}
	\begin{enumerate}
		\item Show that complex conjugation is a linear map $\bC \to \bC$ if we view $\bC$ as a module over $\bR$.
		\item Show that complex conjugation is not linear if we view $\bC$ as a module over $\bC$.
	\end{enumerate}
\end{exercise}
\section{Subspaces and Quotient Spaces}
\section{Bases and Dimension}
\begin{importantdefinition}{Linear Combinations}{}
	Let $M$ be an $R$-module, let $(m_i)_{i \in I} \subseteq M$ and $(r_i)_{i \in I} \subseteq R$. Then a \tib{linear combination} is a sum of the form
	\begin{align*}
		\sum_{i \in I} r_i.m_i
	\end{align*}
\end{importantdefinition}
\begin{importantdefinition}{Vector Span}{}
	Let $M$ be an $R$-module, let $S = (m_i)_{i \in I} \subseteq M$. Then the span of $S$ is the set $\angles{S}$ of all linear combinations over elements of $S$, i.e.
	\begin{align*}
		\angles{S} = \lr{\sum_{i \in J \subseteq I} r_i.m_i \mid (r_i)_{i \in I} \subseteq R, m_i \in S}
	\end{align*}
\end{importantdefinition}
For example, the span of one vector $v \in \bR^3$ is the line through the origin on which $v$ lies. The span of two orthogonal vectors $v,w \in \bR^3$ is the plane through the origin on which $v$ and $w$ lie.
\section{Matrices}
\section{Systems of Linear Equations}
\section{The Method of Least Squares}
\chapter{Affine Spaces}
Affine spaces are structures similar to vector spaces, but where the notions of "points" and "directions" are separate.
\begin{importantdefinition}{Affine Space}{}
	A $\Bbbk$-\tib{affine space} is a tuple $(A, \vA)$, consisting of a set $A$ of \tib{points} and a $\Bbbk$-vector space $\vec A$ of \tib{translation vectors}, alongside a translation map $+ : A \times \vec A \to A$, such that:
	\begin{enumerate}
		\item Translation is associative:
		\begin{align*}
			\forall a \in A, \vv,\vw \in \vA:\\
			(a + \vv) + \vw = a + (\vv + \vw)
		\end{align*}
		\item For any pair of points, there exists a unique vector translating one point onto the other:
		\begin{align*}
			\forall a,b \in P_A:\\
			\exists! \vv \in \vA : a + \vv = b
		\end{align*}
	\end{enumerate}
\end{importantdefinition}
The dimension of an affine space is simply the dimension of its associated vector space. 
\newpar
Any vector space $V$ can be canonically turned into an affine space by taking $V$ as both the set of points and the set of vectors. On the flipside, we can turn any affine space $(A, \vA)$ into a vector space by restricting our attention to translations of one "base point" $b \in A$, and then identifying any point $p \in A$ with the vector $\vv \in \vA$ translating $b$ onto $p$.
\newpar
For most practical purposes, it is therefore sufficient to work with a vector space as your base space for whatever math you are doing, but the non-canonicity of having to choose an arbitrary base point $b \in A$ to turn an affine space into a vector space means that some people prefer affine spaces for philosophical reasons - after all, what is the "base point" of our physical universe?
\newpar
The distinction becomes a lot more important when considering affine subspaces, which are essentialy vector subspaces translated by arbitrary points:
\begin{importantdefinition}{Affine subspaces}{}
	Let $(A, \vA)$ be an affine space. Then an affine subspace of $(A, \vA)$ is a tuple $(B, \vB)$ such that:
	\begin{enumerate}
		\item $B \subseteq A$.
		\item $\vB$ is a vector subspace of $\vA$.
		\item $\vB$ is the set of translations between points in $B$, i.e. for any $b \in B$, $\vv \in \vA$, we have $b + \vv \in B$ if and only if $\vv \in \vB$.
	\end{enumerate}
	Thus, an affine subspace can always be written as
	\begin{align*}
		B = b + \vB := \lr{b + \vv \mid \vv \in \vB \subseteq \vA}
	\end{align*}
	for an arbitrary base point $b \in B$.
\end{importantdefinition}
An affine subspace of dimension $1$ is known as a \tib{line}. An affine subspace of dimension $2$ is known as a \tib{plane}. An affine subspace whose dimension is one lower than the space it resides in is known as an \tib{hyperplane}.
\begin{definition}
	Two affine subspaces $(A, \vA)$, $(B, \vB)$ are called \tib{parallel} if $\vA = \vB$.
\end{definition}
\begin{importantdefinition}{Affine Maps}{}
	Let $A,B$ be $\Bbbk$-affine spaces. Then a map $F : A \to B$ is called \tib{affine} if there exists a linear map $\vF : \vA \to \vB$ such that
	\begin{align*}
		\forall a \in P_a, \vv \in \vA:
		F(a + \vv) = F(a) + \vF(\vv)
	\end{align*}
\end{importantdefinition}
As previously mentioned, if our affine space is a vector space, we can identify any point in $P_A$ with the vector translating the origin onto that point, which lets us write:
\begin{align*}
	F(\vv) := F(0 + \vv) = F(0) + \vF(\vv) := \vF(\vv) + \vw
\end{align*}
thus, an affine map between vector spaces is simply a linear map composed with a translation.
\begin{importanttheorem}{Intercept Theorem}{}
	Let $v_1, v_2 \in \bR$ be linearly independent. Let $L_1$ be an affine line through $v_2$ and $v_1$. Then if we have $\alpha, \beta \in \bR_{\neq 0}$, and $L_2$ is a parallel line through $\alpha v_2$ and $\beta v_1$, we must have $\alpha = \beta$.
\end{importanttheorem}
\begin{proof}
	Since $L_1$ and $L_2$ are parallel, their direction vectors must point in the same direction, which means there must be a $\mu \in \bR_{\neq 0}$ such that
	\begin{align*}
		\mu(v_2 - v_1) &= \mu(\alpha v_2 - \beta v_1)\\
		\implies \mu v_2 - \mu v_1 &= \alpha v_2 - \beta v_1\\
		\implies \mu v_2 - \alpha v_2 &= \mu v_1 - \beta v_1\\
		\implies (\mu - \alpha) v_2 &= (\mu - \beta) v_1
	\end{align*}
	Since $v_2$ and $v_1$ are linearly independent, this implies $\mu- \alpha = \mu - \beta = 0$, i.e. $\alpha = \beta = \mu$.
\end{proof}
\begin{theorem}
	An injective map $\Phi : \bR^2 \to \bR^2$ is linear if and only if it maps $\vzero$ to $\vzero$ and lines to lines.
\end{theorem}
\begin{proof}
	\begin{enumerate}
		\item[$\Rightarrow$]
		Let $\Phi$ be linear. Let $g := \lr{\vp + \lambda \cdot \vv}$ be a line. Then we have
		\begin{align*}
			\Phi(g) 
			&= \lr{\Phi(\vp + \lambda \cdot \vv)}\\
			&= \lr{\Phi(\vp) + \lambda \cdot \Phi(\vv)},
		\end{align*}
		i.e. $\Phi(g)$ is also a line.
		\item[$\Leftarrow$]
		Since $\Phi$ maps $\vzero$ to $\vzero$ and lines to lines, it must map lines through the origin to lines through the origin. Let
		\begin{align*}
			\Phi(e_1) &:= v_1\\
			\Phi[\lambda e_1] &:= L_1 = \lr{\lambda v_1}\\
			\Phi(e_2) &:= v_2\\
			\Phi[\lambda e_2] &:= L_2 = \lr{\lambda v_2}
		\end{align*}
		Since $\Phi$ is injective, $L_1$ and $L_2$ can only intersect at $\lr{0}$, and thus $v_1$ and $v_2$ must be linearly independent.
		\newpar
		Now, let's take a closer look at images of individual points.
		Since $\Phi[\lambda e_1] = \lr{\lambda v_1}$ and $\Phi[\lambda e_2] = \lr{\lambda v_2}$, we must have
		\begin{align*}
			\Phi[\lambda e_1] &= \phi_1(\lambda) v_1\\
			\Phi[\lambda e_2] &= \phi_2(\lambda) v_2
		\end{align*}
		for some $\phi_1, \phi_2 : \bR \to \bR$.
		\newpar
		Now, notice that the line through $\lambda e_1$ and $\lambda e_2$ is parallel to the line through $e_1$ and $e_2$. Since $\Phi$ is injective, it must preserve parallelism, and so the line through $v_1$ and $v_2$ must be parallel to the line through $\phi_1(\lambda)v_1$ and $\phi_2(\lambda)v_2$. The intercept theorem now gives us $\phi_1(\lambda) = \phi_2(\lambda) = \mu$. We thus write $\phi(\lambda) := \phi_1(\lambda) = \phi_2(\lambda)$.
		\newpar
		Since $\Phi$ is injective, $\phi$ must be as well, since otherwise we would have some $v$ with $\Phi(\lambda_1 v) = \Phi(\lambda_2 v)$ despite $\lambda_1 v \neq \lambda_2 v$. We have also have $\phi(0) = 0$ and $\phi(1) = 1$, since $\Phi(0) = 0$ and $\Phi(e_1) = v_1$.
		\newpar
		Now, we want to show that $\Phi(v+w) = \Phi(v) + \Phi(w)$ holds for linearly independent $v,w$. Since $\Phi$ preserves parallelism, it must map parallelograms to parallelograms. Therefore, the parallogram spanned by $v$ and $w$ must be mapped to the parallogram spanned by $\Phi(v)$ and $\Phi(w)$. This means that the upper corner $v+w$ gets mapped to the upper corner $\Phi(v) + \Phi(w)$ which gives us $\Phi(v+w) = \Phi(v) + \Phi(w)$ as desired.
		\newpar
		Putting our insights together, we now know that
		\begin{align*}
			\Phi(\lambda_1 e_1 + \lambda_2 e_2)
			&= \Phi(\lambda_1 e_1) + \Phi(\lambda_2 e_2)\\ 
			&= \phi(\lambda_1) \Phi(e_1) + \phi(\lambda_2) \Phi(e_2)
		\end{align*}
		it only remains to be shown that $\phi$ is the identity. For this, consider that we have:
		\begin{align*}
			(\phi(a) + \phi(b))\Phi(e_1)
			&= \phi(a) \Phi(e_1) + \phi(b) \Phi(e_1)\\
			&= \Phi(ae_1 + be_1\\
			&= \Phi((a+b)e_1)\\ 
			&= \phi(a+b) \Phi(e_1),
		\end{align*}
		which gives us $\phi(a) + \phi(b) = \phi(a+b)$. For $a \neq 0$, the line through $e_1$ and $ae_2$ must point in the same direction as the line through $be_1$ and $bae_2$, so the line through $\Phi(e_1)$ and $\phi(a) \Phi(e_2)$ must be parallel to the line through $\phi(b)\Phi(e_1)$ and $\phi(ab)\Phi(e_2)$, from which the intercept theorem gives us
		\begin{align*}
			\frac{\phi(ab)}{\phi(a)} &= \phi(b)\\
			\implies \phi(ab) &= \phi(a) \phi(b)
		\end{align*}
		Since we have $\phi(0) = 0$, $\phi(1) = 1$, $\phi(a + b) = \phi(a) + \phi(b)$, $\phi(ab) = \phi(a)\phi(b)$,
		our map $\phi : \bR \to \bR$ must be a field homomorphism, of which the identity is the only one. Thus, $\Phi$ is linear.
	\end{enumerate}
\end{proof}
\begin{exercise}
	An injective map between real affine spaces of dimension $\geq 2$ is affine if and only if it maps lines to lines.
\end{exercise}
\chapter{Determinants}
\section{Defining Volume}
\section{Determinants of Endomorphisms}
\section{Orientation}
\chapter{Eigenvectors and Eigenvalues}
\begin{importantdefinition}{Eigenvalues and Eigenvectors}{}
	Let $k$ be a field, $F : V \to V$ an Endomorphism, and $\lambda \in k$. Then the \tib{$\lambda$-eigenspace} of $F$ is the set
	\begin{align*}
		\Eig(F, \lambda) 
		&:= \lr{\vv \in V : F(\vv) = \lambda \vv}\\ 
		&= \lr{\vv \in V : (F - \lambda \id_V)(\vv) = 0}\\
		&= \ker(F - \lambda \id_V)
	\end{align*}
	If $F(\vv) = \lambda \vv$, with $\vv \neq \vzero$, we say that $\vv$ is an \tib{eigenvector} of $F$ with \tib{eigenvalue} $\lambda$.
\end{importantdefinition}
Note that we exclude $\vzero$ from being an eigenvector, since $F(\vzero) = \lambda \vzero$ trivially holds for every endomorphism. Thus every $\lambda$ would be an eigenvalue, which would make the definition pointless. We do, however, allow $0$ as an eigenvalue - the $0$-eigenspace of $F$ is simply its kernel. Sometimes it ends up being useful to rephrase our definition of eigenspaces in terms of kernels:
\begin{align*}
	\Eig(F, \lambda) 
	&= \lr{\vv \in V : F(\vv) = \lambda \vv}\\
	&= \lr{\vv \in V : F(\vv) - \lambda \vv = 0}\\
	&= \lr{\vv \in V : (F - \lambda \id_V)(\vv) = 0}\\
	&= \ker(F - \lambda \id_V)
\end{align*}
This shows that $\Eig(F, \lambda)$ is a subspace of $V$ - by definition, these subspaces must then be the maximal subspaces that are closed under application of $F$, i.e. $F(\Eig(F, \lambda)) = \Eig(F, \lambda)$.
\begin{exercise}
	In an informal, intuition-based way:
	\begin{enumerate}
		\item Find the real eigenvalues and eigenspaces of a reflection of $\bR^2$ along an axis $\lr{\lambda \vv : \lambda \in \bR}$
		\item Find a rotation of $\bR^2$ that has no real eigenvalues. Find two different rotations of $\bR^2$ that do have real eigenvalues. Show that any rotation of $\bR^2$ has complex eigenvalues.
		\item Find the eigenvectors of the derivative map $\frac{d}{dx}$
	\end{enumerate}
\end{exercise}
Our characterisation $\Eig(F, \lambda) = \ker(F - \lambda \id_V)$ means that, if $F$ is finite-dimensional, we can find all eigenvectors of $F$ with a \ul{given} eigenvalue $\lambda$ by solving $(F - \lambda \id_V)\vv = 0$, which is simply a systen of linear equations. Finding the correct eigenvalues ends up being more tricky - in the section after this one, we will work towards proving the following key result:
\begin{proposition}
	The eigenvalues of $F$ are given by the zeroes of the \tib{characteristic polynomial} $\chi_F(\lambda) := \det(F - \lambda \id_V)$.
\end{proposition}
However, before we do so, we want to familiarize ourselves with some of the more basic properties of eigenvectors and eigenvalues.
\newpar
One important application is that finding eigenvectors lets us diagonalize endomorphisms, greatly easing computation afterwards:
\begin{theorem}
	Let $V$ be a vector space over a field $k$. Let $B = (b_1, \hdots, b_n)$ be a basis of $V$, and let $F$ be an endomorphism represented in $B$ by a matrix $A$. Then:
	\begin{enumerate}
		\item A basis vector $b_i$ is an eigenvector of $F$ iff. the $j$-th column of $A$ has the form $\lambda e_j$.
		\item $F$ is diagonalizable iff. there exists a basis of $V$ consisting of eigenvectors of $F$.
	\end{enumerate}
\end{theorem}
\begin{theorem}
	Let $V$ be a vector space over a field $k$. Let $F : V \to V$ be an endomorphism. Then eigenvectors of $F$ with different eigenvalues are linearly independent.
\end{theorem}
\begin{corollary}
	Let $V$ be an $n$-dimensional vector space over a field $k$. Let $F : V \to V$ be an endomorphism. Then:
	\begin{enumerate}
		\item $F$ has no more than $n$ distinct eigenvalues,
		\item If $(\lambda_i)_{i \in I}$ are mutually distinct eigenvalues of $F$, the sum
		\begin{align*}
			\sum_{i \in I} \Eig(F, \lambda_i)
			\subseteq
			V
		\end{align*}
		is direct.
		\item $F$ is diagonalizable if and only if $V$ is a direct sum of eigenspaces of $F$.
		\item If $F$ has $n$ distinct eigenvalues, it is diagonalizable.
	\end{enumerate}
\end{corollary}
\begin{corollary}
	Let $V,W$ be $n$-dimensional vector spaces over a field $k$. Let $F: V \to V$, $G : W \to W$ be endomorphisms. Then there exists a unique endomorphism $P : V \to W$ such that $P \circ F = G \circ P$ if and only if $F$ and $G$ have the same eigenvalues.
\end{corollary}
\section{Characteristic Polynomials}
\begin{importanttheorem}{Characterstic Polynomial}{}
	Let $V$ be a finite-dimensional vector space over a field $\Bbbk$. Let $F : V \to V$ be an endomorphism. Then $\lambda \in \Bbbk$ is an eigenvalue of $F$ if and only if
	\begin{align*}
		\chi_F(\lambda) = \det(\lambda \cdot \id_V - F) = 0_\Bbbk
	\end{align*}
	We call $\chi_F$ the \tib{characteristic polynomial of $F$}.
\end{importanttheorem}
\begin{proof}
	We know that $\lambda \in \Bbbk$ is an Eigenvalue of $F$ if and only if $$\ker(F - \lambda \id_V) \neq \lr{0_V},$$ i.e. if and only if $F - \lambda \id_V$ isn't injective.
	\newpar
	We also know that $F - \lambda \id_V$ isn't injective if and only if it isn't surjective, i.e. if and only if it isn't invertible, i.e. if and only if $$\det(\lambda.\id_V - F) = \chi_F(\lambda) = 0_\Bbbk.$$
\end{proof}
\begin{exercise}
	Show that every endomorphism $F : V \to V$ of a finite-dimensional complex vector space $V$ has at least one eigenvalue.
\end{exercise}
\begin{exercise}
	Show that the constant term of a characteristic polynomial $\chi_F(\lambda)$ is given by $(-1)^n \cdot (\det F)$.
\end{exercise}
(...)
\section{Algebras and the Cayley-Hamilton Theorem}
\begin{importanttheorem}{Evaluation map}{}
	Let $R$ be a commutative unital ring, and let $A$ be a commutative unital algebra over $R$. Then, for every $a \in A$, there exists a unique algebra homomorphism $\ev_a : R[X] \to A$ such that:
	\begin{enumerate}
		\item $\ev_a(r) = 1_A . r$ for all constant polynomials $r \in R[X]$,
		\item $\ev_A(X) = X$.
	\end{enumerate}
	We call $\ev_a$ the \tib{evaluation map of $R[X]$ at $a$}.
\end{importanttheorem}
Informally, this means that we can take elements of an algebra $A$ over $R$ and plug them into polynomials that were originally only defined on $R$, and we won't run into issues. In particular, we can plug matrices into polynomials over an underlying ring! Since $\ev_a$ is unique, we will generally just write $P(a) := \ev_a(P)$ from now on.
\begin{importanttheorem}{Cayley-Hamilton}{}
	Let $R$ be a commutative unital ring. Let $A \in R^{n \times n}$. Then $\chi_A(A) = 0$.
\end{importanttheorem}
\section{Minimal Polynomials}
\begin{importanttheorem}{Minimal polynomial}
	Let $F$ be an endomorphism of a vector space $V$ over a field $\Bbbk$. Then there exists a unique polynomial $\mu_F \in \Bbbk[X]$ such that:
	\begin{enumerate}
		\item $\mu_F$ is either normalized or equal to $0_{\Bbbk[X]}$
		\item $\mu_F$ is a divisor of every $P \in \Bbbk[X]$ with $P(F) = 0_{\End_{\Bbbk}(V)}$
	\end{enumerate} 
	We call $\mu_F$ the \tib{minimal polynomial of $F$}.
\end{importanttheorem}
\chapter{Normal Forms of Endomorphisms}
\section{Decomposing Vector Spaces using Polynomials}
\begin{exercise}
	Let $G$ be an abelian group. Show that $G$ becomes a $\bZ$-module when equipped with the operation $n.g = \underbrace{g + \hdots + g}_{n \tn{ times}}$.
\end{exercise}
\begin{exercise}
	Let $V$ be a vector space over a field $\Bbbk$. Show that $V$ becomes an $\End(V)$-Module when equipped with the operation $F.v = F(v)$.
\end{exercise}
\begin{theorem}
	Any $\Bbbk[X]$-module is the same thing as a $\Bbbk$-vector space $V$ equipped with a specified endomorphism $F$, and vice versa.
\end{theorem}
\begin{proof}
	\begin{enumerate}
		\item Let $V$ be a $\Bbbk$-vector space. We know from the previous exercise that $V$ is an $\End(V)$-module. Thus, the polynomial evaluation map $\ev_F : \Bbbk[X] \to \End(V)$ turns $V$ into $\Bbbk$-module via
		\begin{align*}
			P.v = \ev_F(P)(v) = P(F)(v)
		\end{align*}
		\item Now let $M$ be a $\Bbbk[X]$-module. Then $M$ becomes a $\Bbbk$-module, i.e. a $\Bbbk$-vector space, if we restrict scalar multiplication to only accept constant polynomials. Then it can easily be seen that $G(v) = X.v$ is a linear transformation.
	\end{enumerate}
	\newpar
	Now, note that if we start with a vector space $(V,F)$, then use our construction to get a module $M$, and then turn it back into a vector space, we will obtain our original $F$ again, since $$G(v) = X._Mv = X(F)(v) = F(v)$$ 
	Conversely, if we start with a module $M$, then turn it into a vector space $(V,F)$, and then back into a module $M$, we reobtain our original scalar multiplication, since 
	\begin{align*}
		P._Vv = P(G)(v) = P(X._M)(v) = P._Mv
	\end{align*}
	Therefore, the two directions of our construction are inverse to each other, i.e. we actually have a bijection.
\end{proof}
Thus, given a $\Bbbk[X]$-module $M$, we can write $M = (V,F)$
\begin{exercise}
	Let $M = (V,F)$ and $N = (W,G)$ be $\Bbbk[X]$-Modules. Then a $\Bbbk[X]$-linear map $H : M \to N$ is $\Bbbk$-linear and additionally fulfills $H \circ F = G \circ H$, and any $\Bbbk$-linear map $H$ fulfilling these properties is $\Bbbk[X]$-linear.
\end{exercise}
Let $M = (V,F)$ be a $\Bbbk[X]$-Module. Then, by definition, a subspace $U \subseteq V$ is a $\Bbbk[X]$-submodule of $M$ if and only if $U$ is closed under scalar multiplication with elements of $\Bbbk[X]$.
\begin{exercise}
	Show that $U$ is closed under scalar multiplication with elements of $\Bbbk[X]$ if and only if $Xu \in U$, which itself holds if and only if $F(u) \in U$ holds for all $u \in U$, i.e. $U$ is \tib{$F$-invariant}.
\end{exercise}
\begin{exercise}
	\begin{enumerate}
		\item Show that the scalar multiplication map $P. = P(F)$ is linear.
		\item Conclude that $\ker(P(F))$ and $\im(P(F))$ are submodules of $M$.
		\item Conclude that $\ker(P(F))$ and $\im(P(F))$ are $F$-invariant. 
	\end{enumerate}
\end{exercise}
\begin{importanttheorem}{Structure theorem for modules over a Euclidean ring}{}
	Let $M$ be a module over a euclidean domain $R$. Let $a = bc \in R$ such that $b$ and $c$ share no divisors and $a.m = 0$ for all $m \in M$. Then we have:
	\begin{enumerate}
		\item $\ker(b.) = \im(c.)$,
		\item $\im(b.) = \ker(c.)$,
		\item $M = \ker(b.) \oplus \ker(c.)$.
	\end{enumerate}
\end{importanttheorem}
\begin{proof}
	\begin{enumerate}
		\item Let $c.m \in \im(c.)$ for $m \in M$. Then we have $b.(c.m) = a.m = 0$, i.e. $\im(c.) \subseteq \ker(b.)$. Now, applying the Euclidean algorithm gives us $1 = k_1 b + k_2 c$. If we now let $m \in \ker(b.)$, we get:
		\begin{align*}
			m 
			&= 1m \\
			&= (k_1b + k_2c)m \\
			&= k_1\underbrace{(bm)}_{= 0} + k_2(cm) \\
			&= c(k_2m) \in \im(c.)
		\end{align*}
		Thus, we also have $\ker(b.) \subseteq \im(c.)$, i.e. $\ker(b.) = \im(c.)$.
		\item $\ker(c.) = \im(b.)$ follows from the same argument by swapping the roles of $b$ and $c$.
		\item Now, let $m \in M$. Then, applying our identity from before, we get
		\begin{align*}
			m &= (k_1b + k_2c)m\\
			&= b(k_1m) + c(k_2m)\\
			&\in \im(b.) + \im(c.)\\
			&= \ker(b.) + \ker(c.),
		\end{align*}
		which gives us $M = \ker(b.) + \ker(c.)$ as desired.
		\newpar
		If we have $m \in \ker(b.) \cap \ker(c.)$, we get
		\begin{align*}
			m = k_1(b(m)) + k_2(c(m)) = 0,
		\end{align*}
		so $\ker(b.) \cap \ker(c.) = \lr{0}$, so our sum is direct.
	\end{enumerate}	
\end{proof}
\begin{corollary}
	Let $G$ be an abelian group. Let $n = mk \in \bZ$ such that $m$ and $k$ share no divisors and $n \cdot g = 0$ for all $g \in G$. Then we have
	$$G = \ker(m\cdot) \oplus \ker(k\cdot).$$
\end{corollary}
\begin{importantcorollary}{Chinese Remainder Theorem}{}
	
\end{importantcorollary}
\begin{corollary}
	Let $V$ be a vector space over a field $\Bbbk$ and $F : V \to V$ be an Endomorphism. Let $P = QR \in \Bbbk[X]$ such that $Q$ and $R$ share no divisors and $P.v = P(F)(v) = 0$ for all $v \in V$. Then we have
	$$V = \ker(Q(F)) \oplus \ker(R(F)).$$
\end{corollary}
\begin{importantcorollary}{$P$-Generalized Eigenspace Decomposition}{}
Let $V$ be a vector space over a field $\Bbbk$ and $F : V \to V$ be an Endomorphism. Let $P \in \Bbbk[X]$ be a normed polynomial and let $P = P_1^{n_1} \cdot \hdots \cdot P_k^{n_k}$ be its prime factorization. Then we have
\begin{align*}
	V = \bigoplus_{i = 1}^k \ker(P_i^{n_i}(F))
\end{align*}
\end{importantcorollary}
\section{The Jordan Normal Form}
\begin{importantcorollary}{Generalized Eigenspace Decomposition}{res:geneigenspace}
	\label{res:geneigenspace}
	Let $V$ be a vector space over a field $\Bbbk$ and $F : V \to V$ be an Endomorphism. Assume that its minimal polynomial $\mu_F$ splits into linear factors over $\Bbbk$, i.e. $$\mu_F(x) = \prod_{i = 1}^k(x - \lambda_i)^{n_i}$$ with $\lambda_i \in \Bbbk$. Then we have
	\begin{align*}
		V = \bigoplus_{i = 1}^k \ker((F - \lambda_i \id_V)^{n_i})
	\end{align*}
	We call $\ker((F - \lambda_i \id_V)^{n_i})$ the \tib{generalized $\lambda_i$-Eigenspace of $F$} and denote it $H(F, \lambda_i)$ (from the German "Hauptraum").
\end{importantcorollary}
\begin{exercise}
	Show that $H(F,\lambda_i)$ contains all eigenvectors of $F$ with eigenvalue $\lambda_i$.
\end{exercise}
\begin{exercise}
	\label{res:hauptrauminvariant}
	Show that $H(F,\lambda_i)$ is an $F$-invariant subspace of $V$, i.e. a subspace such that $F(H(F,\lambda_i)) \subseteq H(F,\lambda_i)$
\end{exercise}
\begin{exercise}
	\label{res:matrixdirectsum}
	Let $V = V_1 \oplus V_2$ be a decomposition of $V$ into $F$-invariant subspaces. Let $B_1$ be a basis of $V_1$, and let $B_2$ be a basis of $V_2$. Show that:
	\begin{enumerate}
		\item $B = B_1 \cup B_2$ is a basis of $V$
		\item If the restriction $F|_{V_1}$ of $F$ to $V_1$ is represented in basis $B_1$ as a matrix $A_1$, and the restriction $F|_{V_2}$ of $F$ to $V_2$ is represented in basis $B_2$ as a matrix $A_2$, then $F$ is represented in basis $B$ as
		\begin{align*}
			A = 
			\begin{pmatrix}
				A_1 & 0\\
				0 & A_2
			\end{pmatrix}
		\end{align*}
	\end{enumerate}
\end{exercise}
\begin{exercise}
	\label{res:twobases}
	Let $F : V \to W$ be a linear map. Let $A$ be a basis of $\ker~ F$, and let $B$ be a tuple of vectors such that $F(B)$ is a basis of $\im~ F$. Then $A \cup B$ is a basis of $V$.
\end{exercise}
We now want to make use of this to construct a Jordan basis $B$ of a nilpotent Endomorphism $N$:
\begin{corollary}
	\label{res:jordanbasis}
	Let $N : V \to V$ be a nilpotent endomorphism. Then there exists a basis $B$ of $V$ such that:
	\begin{enumerate}
		\item $B \cup \lr{0}$ is $F$-invariant, i.e. if $b \in B$, then we have $F(b) \in B$ or $F(b) = 0$.
		\item For every $b \in B$, there exists at most one element $b' \in B$ with $F(b') = b$.
	\end{enumerate}
	We call such a basis a \tib{Jordan basis}.
\end{corollary}
\begin{proof}
	We prove this theorem in three steps: We construct a basis $B$, then show it is $F$-invariant, then show that each element of $B$ has at most one preimage in $B$.
	\begin{enumerate}
		\item  
		We construct our basis by induction over $n = \dim V$.
		\newpar
		If $n = 0$, then the empty set is a basis (since a sum over the empty set equals the additive identity, i.e. the zero vector) which (vacuously) fulfills our desired properties.
		\newpar
		Now, since the kernel of $N$ is non-trivial, we have $\dim \im~ N < \dim V$. Therefore, by induction, we can assume we already have a Jordan basis $A$ of $\im~ N$. We split this basis $A$ into two parts:
		\begin{enumerate}
			\item $A_0 := \lr{a \in A \mid N(a) = 0},$
			\item $A_{\neq 0} := \lr{a \in A \mid N(a) \neq 0}$.
		\end{enumerate}
		$A_0$ is by definition a set of linearly independent vectors contained in $\ker A$. By the Steinitz exchange lemma, we can add a set $C$ of Linearly independent vectors such that $A_0 \cup C$ is a basis of $\ker A$.
		\newpar
		$A_\neq 0$ is the set of all $A$ such that $N(A_{\neq 0}) \subseteq A$. We add a set $D$ containing preimages of all the remaining elements of $A$, which must exist since $A \subseteq \im~ N$. By construction, we now have $F(A_{\neq 0} \cup D) = A$.
		\newpar
		Putting everything together, $A_0 \cup C$ is a basis of $\ker~ N$, and $F(A_{\neq 0} \cup D) = A$, which is a basis of $\im~ N$. Therefore, by exercise \ref{res:twobases}, $B = A \cup C \cup D$ is
 		a basis of $V$.
 		\item
 		Now, $B$ must be $F$-invariant since $A$ is $F$-invariant by definition, $C$ is contained in the kernel of $F$ and thus $F(B) = \lr{0}$, and $D$ consists of preimages of $A$, i.e. $F(D) = A \subseteq B$. 
 		\item
 		Finally, the preimages of $A$ in $B$ are contained in either $A_{\neq 0}$ or $D$, which are disjoint by construction.
\end{enumerate}
\end{proof}
\begin{corollary}
	Let $F : V \to V$ be a matrix whose characteristic polynomial has $n = \dim(V)$ roots. Then there exists a basis $B$ of $V$ such that $_BF_B$ is a block diagonal matrix consisting of Jordan blocks.
\end{corollary}
\begin{proof}
	Generalized Eigenspace decomposition \ref{res:geneigenspace} shows us that $V$ is the direct sum of the generalized Eigenspaces of $F$.
	We know by exercise \ref{res:hauptrauminvariant} that the generalized Eigenspaces of $F$ are $F$-invariant. Therefore, exercise \ref{res:matrixdirectsum} tells us that if we find a basis $B_i$ for each generalized Eigenspace $H(F, \lambda_i)$ such that $F|_{H(F, \lambda_i)}$ is represented by a matrix $J_i$ in Jordan normal form, and we combine our bases into a basis $B = \bigcup_{i = 1}^n B_i$ of $V$ as a whole, then the matrix representation of $F$ in basis $B$ will be a block diagonal matrix consisting of these $J_i$, and thus $F$ will also be in Jordan normal form.
	\newpar
	Now, let $F_i$ be the restriction of $(F - \lambda_i \id_V)$ to $H(F, \lambda_i)$.
	By definition of the generalized Eigenspace, we have
	\begin{align*}
		H(F, \lambda_i) = \ker((F - \lambda_i \id_V)^{n_i})
	\end{align*}
	therefore, by construction, $F_i$ is nilpotent. Therefore, by theorem \ref{res:nilpotentjordan}, there exists a basis $B_i$ such that $_{B_i}{F_i}_{B_i}$ is in Jordan normal form, with zeroes on the diagonal. But then, since we have $$F|_{H(F, \lambda_i)} = F_i + \lambda{\id_{H(F, \lambda_i)}}$$
	the matrix representation of $F|_{H(F, \lambda_i)}$ is also in Jordan form, since we simply have to add $\lambda$ to the diagonal of $_{B_i}{F_i}_{B_i}$.
	\newpar
	Therefore, the restriction to $F$ to any of its generalized Eigenspaces is in Jordan normal form in basis $B_i$. As discussed at the beginning of this proof, this means $F$ is in Jordan normal form in basis $B = \bigcup_{i = 1}^n B_i$.
\end{proof}
\begin{exercise}
	Let $F$ be an endomorphism, and let $J$ be its Jordan normal form. Then:
	\begin{enumerate}
		\item The algebraic multiplicity of each eigenvalue of $F$ corresponds to the number of times that eigenvalue appears on the diagonal of $J$.
		\item The geometric multiplicity of any eigenvalue $\lambda$ of $F$ corresponds to the number of $\lambda$-Jordan blocks with that eigenvalue in $J$.
		\item The exponent of $(X-\lambda)$ in the minimal polynomial of $F$ corresponds to the size the of largest $\lambda$-Jordan block in $J$.
	\end{enumerate}
\end{exercise}
\section{The Weierstrass Normal Form}
\chapter{Inner Product Spaces}
\section{Inner Products}
\begin{definition}
	Let $\Bbbk = \bR$ or $\bC$. Let $V$ and $W$ be vector spaces over $\Bbbk$. Then a map $\phi : V \to W$ is called ($\Bbbk$-) \tib{semilinear} or \tib{antilinear}, if we have
	\begin{enumerate}
		\item \tib{Addivity:} $\forall u,v \in V : \phi(u+v = \phi(u) + \phi(v)$
		\item \tib{Antihomogeneity:} $\forall v \in V, \lambda \in \Bbbk : \phi(\lambda.v) = \ol \lambda . \phi(v)$.
	\end{enumerate}
\end{definition}
Notice that for $\Bbbk = \bR$, there is no difference between linearity and antilinearity, since complex conjugation leaves real numbers unchanged.
\begin{definition}
	Let $\Bbbk = \bR$ or $\bC$. Let $V$ be a $\Bbbk$-vector space. Then a map $s : V \times V \to \Bbbk$ is called a \tib{sesquilinear form} if it is antilinear in the first component and linear in the second component,i.e. the map $s_w : v \to s(v,w)$ is antilinear and the map $s_v : w \to s(v,w)$ is linear.
\end{definition}
\begin{exercise}
	Let $A$ be a matrix with entries in $\Bbbk = \bR$ or $\bC$. Show that the map
	\begin{align*}
		s_A(v,w) := v^*Aw
	\end{align*}
	is sesquilinear. (Recall that $v^*$ stands for the conjugated transpose of $v$.)	
\end{exercise}
The latin prefix "sesqui-" means "one and a half" - a sesquilinear is a map that is linear in one component, and "sort of linear", i.e. "one-half-linear" in the other. Some sources require that a sesquilinear form is linear in the first component and antilinear in the second.
\begin{definition}
	A sesquilinear form $s : V \times V \to \Bbbk$ is called a \tib{Hermitian form} if it is \tib{conjugate symmetric}, i.e.
	$s(v,w) = \ol{s(w,v)}$.
\end{definition}
Once again, notice that for $\Bbbk = \bR$, this just means $s(v,w) = s(w,v)$.
\begin{exercise}
	Show that a matrix is Hermitian (i.e. fulfills $A = A^*$) if and only if its associated sesquilinear form $s_A(v,w) = v^*Aw$ is Hermitian.
\end{exercise}
\begin{definition}
	A sesquilinear form $s : V \times V \to \Bbbk$ is called \tib{positive semidefinite} if $s(v,v) \geq 0$ for all $v \in V$. It is called \tib{positive definite} if it additionally fulfills $s(v,v) = 0$ if and only if $v = \vzero$.
\end{definition}
Positive definiteness is the primary reason we are working with sesquilinear forms instead of bilinear forms - we need to be able to compare $s(v,v)$ to $0$, and there is no way to order all of $\bC$ that is compatible with its field structure. However, if we have a sesquilinear Hermitian map, we get $s(v,v) = \ol{s(v,v)}$, i.e. $s(v,v) \in \bR$, which means we can start comparing things again.
\begin{definition}
	A \tib{scalar product} or \tib{inner product} on a vector space $V$ over $\bR$ or $\bC$ is a positive definite Hermitian form. We call a real vector space equipped with a scalar product \tib{Euclidean}, and a complex vector space equipped with a scalar product \tib{unitary}.
\end{definition}
\begin{definition}
	The standard scalar product on $\bC^n$ is given by
	\begin{align*}
		\scalar{v}{w}_\bC = v^* w
	\end{align*}
\end{definition}
\begin{exercise}
	Show that $\scalar{v}{w}_\bC$ is actually a scalar product.
\end{exercise}
\begin{exercise}
	Let $V = C^\infty([0,1],\bR)$ be the room of infinitely differentiable functions $[0,1] \to \bR$. Show that the \tib{$L^2$} scalar product
	\begin{align*}
		\scalar{f}{g}_{L^2} = \int_0^1 f(t) g(t) ~dt
	\end{align*}
	is actually a scalar product.
\end{exercise}
\begin{definition}
	Let $V$ be a vector space over $\Bbbk = \bR$ or $\bC$. Then a \tib{Norm} is a map $\norm{\cdot} : V \to \Bbbk$ fulfilling the following three axioms:
	\begin{enumerate}
		\item \tib{Positive definiteness:} $\norm{v} \geq 0$ and $\norm{v} = 0$ iff. $v = \vzero$
		\item \tib{Homogeneity:} $\norm{\lambda.v} = \abs{\lambda} . \norm{v}$
		\item \tib{Triangle inequality:} $\norm{v + w} \leq \norm{v} + \norm{w}$
	\end{enumerate}
\end{definition}
A norm is a way of assigning lengths to vectors. Once lengths of vectors are defined, one can easily define distances between vectors too:
\begin{exercise}
	Let $\norm{\cdot}$ be a norm on a vector space $V$ over a field $\Bbbk = \bR$ or $\bC$. Show that $d(v,w) := \norm{v - w}$ is a metric on $V$.
\end{exercise}
\begin{importanttheorem}{Cauchy-Schwartz Inequality}{}
	Let $s$ be a scalar product on a vector space $V$ over a field $\Bbbk$. Then we have
	\begin{align*}
		\abs{\scalar{v}{w}} \leq \sqrt{\scalar{v}{v}} \cdot \sqrt{\scalar{w}{w}}
	\end{align*}
\end{importanttheorem}
One important property of scalar products is that they induce norms. Notice that the standard norm on $\bR^n$ is just the square root of the standard scalar product on $\bR^n$. This property generalizes to all scalar products:
\begin{exercise}
	Let $s : V \times V \to \Bbbk$ be a scalar product on a vector space $V$ over a field $\Bbbk$. Let $\norm{v}_s := \sqrt{s(v,v)}$. Show that:
	\begin{enumerate}
		\item $\norm{\cdot}_s$ is a norm on $V$.
		\item We have the \tib{parallelogram identity}
		\begin{align*}
			\norm{v + w}^2_s + \norm{v - w}^2_s = 2\norm{v}_s^2 + 2 \norm{w}_s^2
		\end{align*}
		\item If $\Bbbk = \bR$, we have the \tib{real polarisation formula} \begin{align*}
			s(v,w) = \frac{1}{4}\lr(\norm{v + w}_s^2 - \norm{v - w}_s^2)
		\end{align*}
		\item If $\Bbbk = \bC$, we have the \tib{complex polarisation formula} \begin{align*}
			s(v,w) = \frac{1}{4}\lr(\norm{v + w}_s^2 - \norm{v - w}_s^2
			- i\norm{v + i.w}_s^2 + i\norm{v-i.w}_s^2)
		\end{align*}
	\end{enumerate}
\end{exercise}
\begin{exercise}
	Find a norm on a vector space that cannot be constructed as the square root of a scalar product.
\end{exercise}
\begin{theorem}
	Let $V$ be a vector space over $\Bbbk = \bR$ or $\bC$. Let $\norm{\cdot}$ be a norm on $V$ fulfilling the parallelogram identity. Then there exists a corresponding scalar product $s$ such that $\norm{v} = \sqrt{s(v,v)}$
\end{theorem}
\section{Orthogonality}
\begin{definition}
	Let $(V,s)$ be a vector space with scalar product over a field $\Bbbk = \bR$ or $\bC$. Then we call two vectors $v,w$ \tib{orthogonal} iff.
	\begin{align*}
		s(v,w) = 0
	\end{align*}
\end{definition}
\begin{exercise}
	Let $B = (v_1, \hdots, v_k)$ be a tuple of orthogonal vectors not containing the zero vector. Then $B$ is linearly independent.
\end{exercise}
\begin{definition}
	A basis consisting of pairwise orthogonal vectors is called an \tib{orthogonal basis}. If the basis vectors are also normalized, i.e. have length $1$, we call the basis an \tib{orthonormal basis}.
\end{definition}
\begin{exercise}
	Let $(V,s)$ be a vector space with scalar product.
	Show that $B = (b_1, \hdots, b_n) \subset V$ is an orthonormal basis iff.
	\begin{align*}
		s(b_i,b_j)
		= \delta_{ij}
	\end{align*}
	(Recall that $\delta_{ij}$ is defined as the function that is $1$ if $i = j$ and $0$ otherwise).
\end{exercise}
\begin{definition}
	Let $(V,s)$ be a vector space with scalar product over a field $\Bbbk = \bR$ or $\bC$. Let $U \subseteq V$ be a subspace. Let $B_U := \lr{b_1, \hdots, b_m}$ be a basis of $U$. Then the \tbf{orthogonal projection of $V$ onto $U$} is the map
	\begin{align*}
		p : V &\to U\\
		v &\mapsto \sum_{i = 1}^m g(b_i,v).b_i
	\end{align*}
\end{definition}
\begin{exercise}
	\begin{enumerate}
		\item Show that the orthogonal projection of $V$ onto $U$ leaves vectors in $u$ unchanged.
		\item Show that for each $v \in V$ and $u \in U$, $u$ is orthogonal to $v - p(v)$.
		\item Show that for each $v \in V$, $p(v)$ is the unique point in $U$ for which $\norm{v - p(v)}_s$ is minimal.
	\end{enumerate}
\end{exercise}
\begin{theorem}
	Let $(V,s)$ be a vector space with scalar product over a field $\Bbbk = \bR$ or $\bC$. Let $(v_1, \hdots, v_n)$ be a basis of $V$. Then there are unique vectors $(b_1, \hdots, b_n) \subset V$ such that for all $p \in (1, \hdots, n)$, we have:
	\begin{enumerate}
		\item $(b_1, \hdots, b_p)$ is an $s$-orthonormal Basis of $\angles{v_1, \hdots, v_p} \subseteq V$
		\item $s(b_p,v_p) \in \bR_{> 0}$
	\end{enumerate}
	These vectors $b_p$ can be constructed inductively via $b_1 = \frac{v_1}{\norm{v_1}}$ followed by projection onto the already constructed basis vectors, i.e.
	\begin{align*}
		w_p &= v_p - \sum_{q = 1}^{p-1} g(b_q,v_p).b_q\\´
		b_p &= \frac{w_p}{\norm{w_p}}
	\end{align*}
\end{theorem}
\begin{corollary}
	Every vector space with scalar product has an orthonormal basis.
\end{corollary}
\section{Matrix Representations of Inner Products}
\begin{definition}
	Let $s$ be a sesquilinear form on a finite dimensional vector space $V$over a field $\Bbbk = \bR$ or $\bC$. Let $B = (b_1, \hdots, b_n)$ be a basis of $V$. Then the \tib{Gram matrix of $B$ by $g$} is defined as the matrix with entries
	\begin{align*}
		G_{ij} = s(b_i,b_j) \in \Bbbk^{n \times n}
	\end{align*}
\end{definition}
\begin{exercise}
	Show that the Gram matrix of the standard basis of $\bR^n$ by the scalar product is the identity matrix.
\end{exercise}
\begin{definition}
	A matrix $A \in \Bbbk^{n \times n}$ is called \tib{positive (semi-)definite} if $s_A(v,w) = v^*Aw$ is positive (semi-)definite.
\end{definition}
\begin{exercise}
	Let $s$ be a sesquilinear form on a finite-dimensional vector space $V$ over a field $\Bbbk = \bR$ or $\bC$. Let $G$ be its Gram matrix. Then we have $s(v,w) = v^* G y$ in basis $B$. In particular, $s$ is a scalar product if and only if its Gram matrix is Hermitian and positive definite.
\end{exercise}
\begin{theorem}
	Let $A \in \Bbbk^{n \times n}$, where $\Bbbk = \bR$ or $\bC$. Then the following are equivalent:
	\begin{enumerate}
		\item $A$ is hermitian and psitive definite.
		\item \tib{Cholesky Decomposition:} There exists an upper triangular Matrix $B$ with positive real diagonal entries such that $A = B^*B$.
		\item There exists an invertible Matrix $B$ such that $A = B^* B$.
		\item \tib{Sylvester-Hurwitz Criterion:} $A$ is Hermitian and and all of its square upper left submatrices $A_r$, known as its \tib{leading principal minors}, have strictly positive determinant.
	\end{enumerate}
\end{theorem}
\section{Dual Spaces}
\begin{importantdefinition}{Dual Space}{}
	Let $V$ be a vector space over a field $\Bbbk$. Then we define $V^* = \Hom_{\Bbbk}(V, \Bbbk)$. Elements of $V^*$, i.e. linear maps $V \to \Bbbk$, are known as \tib{linear forms}.
\end{importantdefinition}
\begin{importanttheorem}{Dual Basis}
	Given a basis $B = (b_1, \hdots, b_n)$ of a vector space $V$, the set of linear forms $E = (\epsilon_1, \hdots, \epsilon_n)$ given by
	\begin{align*}
		\epsilon_p(e_q) = \delta_{pq}
	\end{align*}
	is a basis of $V^*$, known as the \tib{dual basis} of $B$.
\end{importanttheorem}
\begin{exercise}
	Let $C^n(\Bbbk, \Bbbk)$, $n \in 0, 1, \hdots, \infty$ be the set of $n$-times continuously differentiable functions $\Bbbk \to \Bbbk$. Then for any $x_0 \in \Bbbk$, the map
	\begin{align*}
		\cdot(x_0): f \to f^{(n)}(x_0)
	\end{align*}
	is a linear form on $C^n(\Bbbk, \Bbbk)$ (where $f^{(n)}$ denotes the $n$-th derivative of $f$).
\end{exercise}
\begin{exercise}
	The map mapping an integrable function $f$ to $\int_{0}^1 f(x) ~dx$ is a linear form. More generally, for any integrable $g : [0,1] \to \Bbbk$, the map taking $f$ to $\int_0^1 \ol{g(x)} \cdot f(x) ~dx$ is a linear form.
\end{exercise}
\begin{theorem}
	Let $(V,g)$ be a $\Bbbk$-vector space with scalar product. Then $g$ induces an injective antilinear map $g_1 : V \to V^*$ via
	\begin{align*}
		g_1(v) = g(v, \cdot)
	\end{align*}
	A linear form $\alpha \in V^*$ is called \tib{representable} by $g$ if there exists a $v \in V$ such that $\alpha = g(v, \cdot)$.
\end{theorem}
\begin{definition}
	Let $V$ be a $\Bbbk$-vector space. Then a \tib{antilinear form} is an antilinear forms $\gamma: V \to \Bbbk$. The \tib{antidual space} $\ol V^*$ of $V$  is the space of antilinear maps $V \to \Bbbk$.
\end{definition}
\begin{theorem}
	Let $(V,g)$ be a finite-dimensional vector space over a field $\Bbbk$. Then the maps $g : V \to V^*$ and $\ol g : V \to V^*$ are bijective.
\end{theorem}
\begin{corollary}
	Let $(V,g)$ be a finite-dimensional vector space. Then all linear and antilinear forms on $V$ are representable by $g$.
\end{corollary}
\begin{definition}
	Let $(V,g)$ and $(W,h)$ be $\Bbbk$-vector spaces with scalar products. Then a linear map $F : V \to W$ is called \tib{adjointable} if there exists a map $G : W \to V$ such that
	\begin{align*}
		g(G(w),v) = h(w,F(v)
	\end{align*}
	for all $v \in V$ and $w \in W$. Such a map $G$ is known as the \tib{adjoint of $F$}, and is written $G = F^*$. If $F = F^*$, we call $F$ \tib{self-adjoint}.
\end{definition}
\begin{proposition}
	If it exists, the adjoint of a given map is unique.
\end{proposition}
\begin{proposition}
	Adjoint maps are linear.
\end{proposition}
\begin{proposition}
	If $G$ is the adjoint of $F$, then $F$ is the adjoint of $G$, i.e. we have $(F^*)^* = F$.
\end{proposition}
\begin{proposition}
	If $A$ is a matrix representation of $F$ by two orthonormal bases, then the conjugate transpose $A^*$ is a matrix representation of $F^*$.
\end{proposition}
\begin{definition}
	Let $V$ and $W$ be $\Bbbk$-vector spaces, with $\Bbbk$ a field. Let $F : V \to W$ be a linear map. Then the \tib{dual} of $F$ is the map 
	\begin{align*}
		F^\vee: W^* &\to V^*\\
		\beta &\mapsto \beta \circ F
	\end{align*}
	i.e.
	\begin{align*}
		F^\vee(\beta)(\vv) = \beta(F(\vv))
	\end{align*}
	for $\beta \in W^*$, $\vv \in V$.
\end{definition}
We can construct a dual of any linear map this way, while linear maps involving infinite-dimensional vector spaces are generally not adjoinable.
\begin{exercise}
	Verify the following properties of the dual map:
	\begin{enumerate}
		\item $(\id_V)^\vee = \id_{V^*}$
		\item Let $F : V \to W$ and $G : U \to V$ be linear, then $(F \circ G)^\vee = G^\vee \circ F^\vee$
		\item The dual of a linear map is linear.
		\item If $F : V \to W$ is given by $F(\vv) = A \cdot \vv$ for some matrix $A$, then $F^\vee(\beta) = \beta \cdot A$, where elements $\beta \in W^*$ are represented as row vectors.
	\end{enumerate}
\end{exercise}
\begin{theorem}
	Let $U \subseteq V$ be vector spaces. Let $\iota : U \to V$ be the inclusion map. Let $q : V \to V / U$ be the quotient map. Then $(V / U)^* \cong \im~ q^\vee = \ker~ \iota^\vee$.
\end{theorem}
\begin{proof}
	Since $q$ maps $V$ to $V/U$, we have $q^\vee : (V/U)^* \to \im~ q^* \subseteq V^*$. Since $q$ is linear, so is $q^\vee$. By definition, any map surjectively maps onto its image. Therefore, all we need to show to obtain $(V / U)^* \cong \im~ q^\vee$ is injectivity.
	\newpar
	So, let $\beta \in \ker~ q^\vee \subseteq (V/U)^*$, i.e. let $q^\vee(\beta) = \beta(q)$ be the zero map. Since $q$ is surjective, $\beta(q(v)) = 0$ for all $v \in V$ already implies $\beta(w) = 0$ for all $w \in V/U$, so $\beta$ must have already been the zero map. 
	\newpar
	Finally, we want to show $\im~ q^\vee = \ker~ \iota^\vee$. 
	\newpar
	First off, recall that $\iota^\vee : V^* \to U^*$. By the universal property of the quotient map, for every $\beta \in \ker \iota^\vee \subseteq V^*$, there exists $\ol \beta : V/U \to \Bbbk$ (i.e. $\ol \beta \in (V/U)^*$) with $\beta = \ol \beta \circ q = \ol \beta(q) \in \im~ q^\vee$, so we have $\ker \iota^\vee \subseteq \im~ q^\vee$.
	\newpar
	Conversely, for any $\beta \in (V/U)^*$ and $v \in (V/U)$, we have
	\begin{align*}
		q^\vee(\beta)(v) = \beta(q(v)) = \beta(0) = 0,
	\end{align*}
	so we also have $\im q^\vee \subseteq \ker \iota^\vee$.
\end{proof}
\begin{importanttheorem}{}{}
	Let $(V,g)$ and $(W,h)$ be $\Bbbk$-vector spaces with scalar products. Let $h_1 : W \to W^*$ and $g_1 : V \to V^*$ be the antilinear maps induced by $g$ and $h$, i.e.
	\begin{align*}
		h_1(\vw) = h(\vw, \cdot)\\
		g_1(\vv) = g(\vv, \cdot)
	\end{align*} 
	Let $F : V \to W$ be linear. Then $g_1 \circ F^* = F^\vee \circ h_1$.
\end{importanttheorem}
\begin{proof}
	Let $v \in V$, $w \in W$, then we have:
	\begin{align*}
		g_1(F^*(w))(v) 
		&= g(F^*(w),v)\\
		&= h(w,F(v))\\
		&= h_1(w)(F(v))\\
		&= F^\vee(h_1(w))(v)
	\end{align*}
\end{proof}
\begin{corollary}
	If $V$ is finite-dimensional, $F : V \to W$ is adjointable.
\end{corollary}
\begin{proof}
	If $V$ is finite-dimensional, $g_1$ is invertible. We get $F^* = g_1^{-1} \circ F^\vee \circ h_1$.
\end{proof}
\begin{definition}
	Let $V$ and $W$ be $\Bbbk$-vector spaces. Let $F : V \to W$ be linear, and let $s : W \times W \to \Bbbk$ be a sesquilinear form. Then we define the \tib{pullback of $s$ by $F$} to be the map $F^*s : V \times V \to \Bbbk$ given by
	\begin{align*}
		(F^*s)(u,v) = s(F(u), F(v))
	\end{align*}
\end{definition}
\begin{exercise}
	Verify that:
	\begin{enumerate}
		\item $F^*s$ is a sesquilinear form.
		\item If $s$ is Hermitian, $F^*s$ is Hermitian.
		\item If $s$ is positive semidefinite, $F^*s$ is positive semidefinite.
		\item 
		If $s$ is positive definite and $F$ is injective, $F^*s$ positive definite.
		\item If we fix bases of $V$ and $W$ such that $F$ is represented by a matrix $A$ and $s$ is represented by a Gram matrix $G$, then $F^*s$ is represented by $A^*GA$.
		\item Given $F : V \to W$ antilinear, we can achieve the same results by defining $F^*s(u,v) = s(F(v),F(u))$.
	\end{enumerate}
\end{exercise}
\chapter{Normal Endomorphisms}
\begin{importantdefinition}{}{}
	Let $(V,g)$ be a $\Bbbk$-vector space with scalar product. Let $F : V \to V$ be an adjointable endomorphism. Then we call $F$ \tib{normal} if $F^* \circ F = F \circ F^*$.
\end{importantdefinition}
\begin{exercise}
	\begin{enumerate}
		\item Show that self-adjoint maps are normal.
		\item A linear map with $F^* = -F$ is called \tib{skew}. Show that skew maps are normal.
		\item A linear map represented in a given basis by a matric of the form
		\begin{align*}
			\begin{pmatrix}
				a & -b\\
				b & a
			\end{pmatrix}
		\end{align*}
		is normal.
	\end{enumerate}
\end{exercise}
\begin{importantdefinition}{}{}
		Let $(V,g)$ and $(W,h)$ be $\Bbbk$-vector spaces with scalar products. Then a linear map $F : V \to W$ is called an \tib{isometric embedding} if $g(v,w) = h(F(v), F(w))$ for all $v,w \in V$.
\end{importantdefinition}
\begin{exercise}
	Show that:
	\begin{enumerate}
		\item Isometric embeddings preserve lengths and distances.
		\item Isometric embeddings preserve angles.
		\item Isometric embeddings are injective.
		\item If orthonormal bases are chosen for $V$ and $W$, then the matrix representation $A$ of $F$ fulfills $A^*A = E_n$.
		\item If $\dim V = \dim W$, $F$ is invertible with $F^{-1} = F^*$. An invertible isometric embedding is called an \tib{isometry}.
		\item Isometries are normal.
	\end{enumerate}
\end{exercise}
\begin{importanttheorem}{Fundamental Theorem of Normal Endomorphisms}{}
	Let $(V,g)$ be a finite-dimensional complex vector space with scalar product. Let $F : V \to V$ be an endomorphism. Then there exists an orthonormal basis of $(V,g)$ consisting of eigenvectors of $F$ if and only if $F$ is normal.
\end{importanttheorem}
Note that this only works for complex numbers, not for reals:
\begin{exercise}
	Recall that any endomorphism represented by a matrix of the form
	\begin{align*}
		A
		=
		\begin{pmatrix}
			a & -b\\
			b & a
		\end{pmatrix}
	\end{align*}
	is normal. Show that, if $b \neq 0$, $A$ has no real eigenvalues.
\end{exercise}
\begin{importantcorollary}{A Normal Form for Complex Normal Endomorphisms}{}
	Let $(V,g)$ be a finite-dimensional real vector space with scalar product. Let $F : V \to V$ be an endomorphism. Then there exists an orthonormal basis of $(V,g)$ in which $F$ is represented by a block diagonal matrix consisting of $1 \times 1$ blocks and $2 \times 2$ blocks of the form
	\begin{align*}
		\begin{pmatrix}
			a & -b\\
			b & a
		\end{pmatrix}
	\end{align*}
	with $b > 0$ if and only if $F$ is normal. If such a representation exists, it is unique up to permutation of the blocks.
\end{importantcorollary}
\begin{importantcorollary}{Spectral Theorem for Self-Adjoint Maps}{}
	Let $(V,g)$ be a finite dimensional real or complex vector space with scalar product. Let $F : V \to V$ be an endomorphism. Then there exists an orthonormal basis of $V$ in which $F$ is represented by a diagonal matrix with real entries if and only if $F$ is self-adjoint.
\end{importantcorollary}
\begin{exercise}
	Derive an analogous spectral theorem for skew maps.
\end{exercise}
\begin{corollary}
	Let $(V,g)$ be a finite-dimensional real or complex vector space with scalar product. Let $s$ be a sesquilinear form on $V$. Then there exists a $g$-orthonormal Basis of $V$ in which $s$ is represented by a diagonal matrix with real entries if and only $s$ is hermitian.
\end{corollary}
\begin{importantcorollary}{Singular Value Decomposition}{}
	Let $(V,g)$ and $(W,h)$ be finite-dimensional real or complex vector spaces with scalar products. Let $F : V \to W$ be linear. Then there exist orthonormal bases of $V$ and $W$ in which $F$ is represented by a diagonal matrix with uniquely determined real entries $a_n$ such that $a_1 \geq a_2 \geq \hdots$ and $a_n = 0$ if $n \geq \tn{rk}~F$. These entries are known as the \tib{singular values} of $F$.
\end{importantcorollary}
Singular value decomposition is fundamental to the JPG compression algorithm, which represents an image by a matrix, then calculates the singular value decomposition of that matrix, and then sets all of the singular values below a given threshold to zero, leading to less data needing to be stored while still maintaining the singular values that 'noticably contribute' to the resulting image.
\section{Linear Isometries}
\begin{importantcorollary}{Normal Form for Real Isometries}{}
	Let $(V,g)$ be a finite-dimensional real vector space with scalar product $g$. Let $F : V \to V$ be a linear isometry. Then there exists an orthonormal basis of $(V,g)$ in which $F$ is represented by a block diagonal matrix consisting of $1 \times 1$ blocks and $2 \times 2$ blocks of the form
	\begin{align*}
		\begin{pmatrix}
			a & -b\\
			b & a
		\end{pmatrix}
	\end{align*}
	with $a^2 + b^2 = 1$.
\end{importantcorollary}
\begin{importantcorollary}{Normal Form for Complex Isometries}{}
	Let $(V,g)$ be a finite-dimensional complex vector space with scalar product $g$. Let $F : V \to V$ be a linear isometry. Then there exists an orthonormal basis of $(V,g)$ in which $F$ is represented by a diagonal matrix whose entries have absolute value $1$.
\end{importantcorollary}
\begin{importantcorollary}{}{}
	Let $(V,g)$ be a finite-dimensional complex vector space with scalar product $g$. Let $F : V \to V$ be a linear isometry. Then $V$ can be decomposed into a direct sum of subspaces $U$ of dimension one or two, such that $F$ acts on the one-dimensional subspaces as the identity or reflection, and on the two-dimensional subspaces as a rotation around the origin.
\end{importantcorollary}
\begin{exercise}
	A linear isometry is orientation-preserving if and only if the number of invariant subspaces on which it acts as a reflection is even.
\end{exercise}
\section{Affine Isometries}
\begin{theorem}
	Let $(A, \vA, s))$, $(B, \vB, g)$ be finite-dimensional normed affine spaces of the same dimension over vector spaces with scalar products. Then there exists an affine isometry $F : A \to B$.
\end{theorem}
\begin{theorem}
	Let $(A, \vA, g)$ be a finite-dimensional normed $\Bbbk$-affine space over a vector space with scalar product. Let $F : A \to A$ be an affine sometry. Then there exists a point $o \in A$ and a vector $x \in \Eig(\vF, 1)$ such that
	\begin{align*}
		F(a) = o + x + \vF(a - o)
	\end{align*}
	$x$ is uniquely determined by $A$, while $o$ is an arbitrary point chosen from an affine-$F$-invariant subspace $(B, \Eig(\vF, 1)) \subseteq (A, \vA)$. 
\end{theorem}
\begin{definition}
	We call $(B, \Eig(\vF, 1))$ the \tib{axis set} of $F$, and call lines of the form
	\begin{align*}
		\lr{b + \lambda.x \mid \lambda \in \Bbbk} \subseteq B
	\end{align*}
	with $b \in B$ the
	\tib{axes} of $F$.
\end{definition}
Our theorem thus states that any affine isometry can be represented as a linear map composed with a translation along the axis set of $F$. 
\typeout{(TODO: Cleaner proof using least squares)}
\begin{proof}
	We write $F(a) = \vec F(a) + y$. Let $U = \Eig(\vec F, 1$). Then $U$ and $U^\perp$ are both invariant under $\vec F$, and $\vec F|_U = \id_U$.
	\newpar
	We write $y = (x,z) \in U \oplus U^\perp$. The map $\id_{U^\perp} - \vec F|_{U^\perp}$ is invertible, since $U^\perp$ doesn't contain any eigenvectors of $\vec F$. We can thus find a $q \in U^\perp$ with $q - \vec F(q) = z$.
	\newpar
	We can then choose any $p \in U$ and set $o = (p,q)$. For any other $v = (u,w) \in V$, we then get
	\begin{align*}
		F(a) 
		&:= F(o + v)\\
		&= y + \vec F(o+v)\\
		&= (x,z) + \vec F((p,q) + (u,w))\\
		&= (x,z) + \vec F((p+u,q+w))\\
		&= (x + \vec F(p) + \vec F(u), z + \vec F(q) + \vec F(w))\\
		&= (x + p + u, z + \vec F(q) + \vec F(w))\\
		&= (x + p + u, (q - \vec F(q)) + \vec F(q) + \vec F(w))\\
		&= (x + p + u, q + \vec F(w))\\
		&= (p,q) + (x,0) + \vec F(u,w)\\
		&= o + (x,0) + \vec F(v),\\
		&= o + x + \vec F(a - o),
	\end{align*}
	which is our desired representation.
	\newpar
	Now, assume we chose an alternate origin $o' = o + (p', q')$ such that
	$$F(a) := F(o' + v) = o' + x' + \vec F(v)$$
	for $v = (u,w)$.
	Then we get:
	\begin{align*}
		x' 
		&= F(o' + v) - o' - \vec F(v)\\
		&= F(o + (p',q') + v) - o - (p',q') - \vec F(v)\\
		&= (o + x + \vec F((p', q') + v)) - o - (p',q') - \vec F(v)\\
		&= (o + x + \vec F(p', q') + \vec F(v) - o - (p',q') - \vec F(v)\\
		&= x + \vec F(p', q') - (p',q')\\
		&= x + \vec (p', \vec F(q')) - (p',q')\\
		&= (x, \vec F(q') - q')
	\end{align*}
	Now, applying $\vec F$ to $x$ and $x'$, we get:
	\begin{align*}
		\vec F(x)
		&=
		F(x) - (x,z)\\
		&=
		F((x,0)) - (x,z)\\
		&=
		o + (x,0) + \vec F((x,0) - o) - (x,z)\\
		&=
		o + (x,0) + (x,0) - \vec F(o) - (x,z)\\
		&=
		o + (x,0) + (x,0) - (p,q-z) - (x,z)\\
		&=
		(p,q) + (x,0) + (x,0) - (p,q) + (0,z) - (x,z)\\
		&=
		(x,0)
		=
		x
	\end{align*}
	and thus
	\begin{align*}
		\vec F(x') &= (\vec F(x), \vec F(\vec F(q')) - \vec F(q'))\\
		&= (x, \vec F(\vec F(q')) - \vec F(q'))
	\end{align*}
	therefore, we get $\vec F(x') = x'$ if and only if $$\vec F(\vec F(q')) - \vec F(q') = \vec F(q') - q',$$
	i.e.
	\begin{align*}
		\vec F(\vec F(q')) - 2\vec F(q') + q' = 0,
	\end{align*}
	i.e.
	\begin{align*}
		(\id_W - \vec F|_W)(\id_W - \vec F|_W(q')) = 0
	\end{align*}
	Since $\id_W - \vec F|_W$ is invertible by construction, this holds iff. $q' = 0$, i.e. if $x' = (x, \vec F(q') - q') = x$ and $o' = o + (p', 0) = o + \lambda x \in B$.
\end{proof}
\chapter{Revisiting Synthetic Geometry}
\section{Incidence Geometries}
\begin{importantdefinition}{Incidence Geometries}{}
	An \tib{incidence geometry} is a set $X$ equipped with a set $G \subseteq \cP(X)$ of subsets of $X$ known as \tib{lines}, such that:
	\begin{enumerate}
		\item Each line contains at least two points,
		\item For every pair of points, there exists exactly one line containing both.
	\end{enumerate}
\end{importantdefinition}
Instead of "$x$ is contained in $g \in G$", we often say "$x$ lies on $g \in G$".
\begin{exercise}
	Let $R$ be a division ring. Then $X := R^n$ can be turned into an incidence geometry by defining lines to be all sets of the form $$g = \lr{p + \lambda v : \lambda \in R},$$ where $p,v \in R^n$.
\end{exercise}
Moving forward, this will be the canonical choice of incidence geometry on $R^n$.
\begin{definition}
	A set of points lying on the same line is called \tib{colinear}.
\end{definition}
\begin{definition}
	A set of lines containing a common point is called \tib{confluent}.
\end{definition}
\begin{definition}
	A set of three non-colinear points is known as a \tib{triangle}.
\end{definition}
\begin{importantdefinition}{Betweenness}{}
	A \tib{betweenness relation} on an incidence geometry $(X,G)$ is a subset $Z \subseteq X^3$ such that the following properties hold:
	\begin{enumerate}
		\item \tib{Induced Orderings on Lines:} On every line $g \in G$, there exists an ordering $\leq$ such that $(x,y,z) \in Z$ if and only if $x \leq y \leq z$ or $x \geq y \geq z$.
		\item \tib{Pasch's Axiom:} Let $p,q,r \in X$. Let 
		$$[p,q] := \lr{x \in X : (p,x,q) \in Z}$$ denote the \tib{line segment between $p$ and $q$}. Then every line $g \in G$ either intersects none of the line segments $[p,q], [q,r], [r,p]$, or at least two.
	\end{enumerate}
	An incidence geometry equipped with a betweenness relation is known as an \tib{ordered incidence geometry}.
\end{importantdefinition}
\begin{exercise}
	Let $K$ be an ordered field. Then:
	\begin{enumerate}
		\item The set of all tuples $(x,y,z)$, where $x,y,z \in K^n$ fulfill $y = \lambda x + \mu z$ with $\lambda, \mu \in [0,1] \subset K$ and $\lambda + \mu = 1$, gives a betweenness relation on $K^n$.
		\item There exists an affine isomorphism $\phi : K \cong (g \in G)$ between any line in $G$ and $K$.
		\item The ordering on $K$ can be recovered from the betweenness relation on $K^n$ by taking the unique ordering $\leq$ on $g \in G$ that fulfills $\phi(0) \leq \phi(1)$.
	\end{enumerate}
\end{exercise}
Moving forward, this will be the canonical choice of betweenness relation on $K^n$.
\begin{definition}
	Given an ordered incidence geometry $(X,G,Z)$, a \tib{ray} in $X$ is a set $A \subseteq X$ of the form
	\begin{align*}
		A = \lr{x \in g \mid x \leq p},
	\end{align*}
	where $g \in G$ is a line and $p \in g$ is a point on that line.
\end{definition}
\begin{importantdefinition}{Congruence Groups}{}
	Given an ordered incidence geometry $(X,G,Z)$, a \tib{congruence group} of $(X,G,Z)$ is a group $K$ of automorphism of $(X,G,Z)$ such that every pair of rays in $A,B \subseteq X$ has exactly two elements $h,k \in K$ such that $h(A) = k(A) = B$.
\end{importantdefinition}
Intuitively, one of these two elements will be a rotation and the other a rotation composed with reflection along an axis.
\begin{exercise}
	Let $K$ be an ordered field. Let $s $ be the standard scalar product in $K^2$. Then the set of $s$-invariant affine automorphisms of $K^2$ forms a congruence group of $K^2$.
\end{exercise}
Moving forward, this will be the canonical choice of congruence group on $K^2$.
\begin{definition}
	We say that an ordered incidence geometry has the \tib{least upper bound property} if every nonempty subset $A \subset g \in G$ that is bounded from above has a least upper bound.
\end{definition}
\begin{definition}
	We call an ordered incidence geometry $(X,G,Z)$ with congruence group $K$ \tib{near-euclidean} if $(X,G,Z)$ fulfills the supremum property and $G$ is nonempty.
\end{definition}
\begin{definition}
	We say that an incidence geometry $(X,G)$ fulfills the \tib{parallel postulate} iff. for every line $g \in G$ and every point $p \notin G$, there exists at most one line $h \in G$ such that $p \in h$ and $h \cap g = \emptyset$.
\end{definition}
\begin{exercise}
	Let $K \subsetneq \bR$ be a proper subfield of $\bR$ containing square roots of every positive element. Then $K^2$ fulfills the parallel postulate.
\end{exercise}
The primary goal of this chapter will be proving the following result:
\begin{importanttheorem}{Classification of near-euclidean geometries}{}
	Let $(X,G,Z)$ be a near-euclidean incidence geometry. Then:
	\begin{enumerate}
		\item If $(X,G,Z)$ fulfills the parallel postulate, it is isomorphic to $\bR^2$.
		\item If $(X,G,Z)$ doesn't fulfill the perallel postulate, it is isomorphic to the \tib{hyperbolic plane} $H$.
	\end{enumerate}
\end{importanttheorem}
\section{Projection Maps}
\section{Projectivization}
\chapter{Hyperbolic Geometry}
\section{Circle Inversions}
\begin{definition}
	Let $E$ be a Euclidean plane. Then the corresponding \tib{extended Euclidean plane} is the set
	\begin{align*}
		\hat E = E \cup \lr{\infty},
	\end{align*}
	adding a "point at infinity" to $E$.
\end{definition}
\begin{definition}
	We call a subset $L \subseteq \hat E$ a \tib{generalized circle} if it is either a circle in $E$ or the union of a line in $E$ with $\lr{\infty}$. We call the former \tib{real circles} and the latter \tib{extended lines}.
\end{definition}
\begin{definition}
	Let $E$ be a euclidean plane. Then we assign to each generalized circle $L \subseteq \hat E$ the \tib{circle inversion map}
	\begin{align*}
		s_L : \hat E \to \hat E,
	\end{align*}
	which is given by the ordinary reflection at $L$ with the additional rule $s_L(\infty) = \infty$ if $L$ is an extended line, and by
	\begin{align*}
		c + \vv \mapsto
		\begin{cases} 
			c + \lr(\frac{r}{\norm{\vv}})^2 \vv & \vv \neq \vzero\\
			\infty &\vv = \vzero
		\end{cases}
	\end{align*}
	with $\infty \mapsto c$ if $L$ is a circle with center $c$ and radius $r$.
\end{definition}
\begin{exercise}
	$s_L$ is an involution, i.e. we have $s_L^2 = \id$.
\end{exercise}
\begin{exercise}
	For any two points $p,q \in \hat E$, there exists a circle inversion $s_L$ with $s_L(p) = q$ and $s_L(q) = p$
\end{exercise}
\begin{exercise}
	Every circle inversion maps generalized circles to generalized circles.
\end{exercise}

\pagebreak
\section{Möbius Transformations}
\begin{definition}
	A composition of circle inversions is known as a \tib{Möbius transformation}.
\end{definition}
\begin{exercise}
	The set $\tn{Möb}(\hat E)$ of Möbius transformations of a given extended Euclidean plane $\hat E$ forms a group under composition.
\end{exercise}
\begin{exercise}
	Any two generalized circles in a generalized Euclidean plane can be mapped onto each other by a Möbius transformation.
\end{exercise}
\begin{theorem}
	Given an extended Euclidean plane $E$, its Möbius transformations are precisely the bijections $\phi : \hat E \iso \hat E$ mapping generalized circles to generalized circles.
\end{theorem}
\begin{theorem}
	Any Möbius transformation that leaves a given generalized circle $L$ pointwise unchanged is either the circle inversion $s_L$ at that circle or the identity map.
\end{theorem}
\begin{exercise}
	Given a generalized Euclidean plane $\hat E$, a Möbius transformation $\phi : \hat E \to \hat E$ and a generalized circle $L \subseteq \hat E$, we have $\phi \circ s_L \circ \phi^{-1} = s_{\phi(L)}$
\end{exercise}
\begin{exercise}
	Given a generalized Euclidean plane $\hat E$, a Möbius transformation $\phi : \hat E \to \hat E$, and a generalized circle $L \subseteq \hat E$, we have $\phi(L) = L$ if and only if $\phi \circ s_L = s_L \circ \phi$.
\end{exercise}
\begin{definition}
	Let $\hat E$ be a generalized Euclidean plane. Given a pair of generalized circles $L, M \subset \hat E$ we call $L$ \tib{orthogonal to $M$}, which we write $L \perp M$, if $L \neq M$ and $s_L(M) = M$.
\end{definition}
\begin{theorem}
	Let $L,M \subset \hat E$ be generalized circles in a generalized Euclidean plane $\hat E$ such that $L \perp M$. Let $\phi : \hat E \to \hat E$ be a Möbius transformation. Then $\phi(L) \perp \phi(M)$.
\end{theorem}
\begin{proposition}
	Let $E$ be a Euclidean plane, $g \subset E$ a line, $H \subset E \setminus g$ an open Half-plane with $g$ as its boundary, and $$K_H = \tn{Möb}_H(\hat E) := \lr{\phi \in \tn{Möb}(\hat E) \mid \phi(H) = H}$$ the set of Möbius transformations leaving $H$ invariant. Then:
	\begin{enumerate}
		\item Any Möbius transformation that leaves $H$ invariant also leaves the extended line $\hat g$ invariant.
		\item $\tn{Möb}_H(\hat E)$ operates transitively on $H$, i.e. any point of $H$ can be mapped onto any other by a Möbius transformation $\phi \in \tn{Möb}_H(\hat E)$.
		\item $\tn{Möb}_H(\hat E)$ is generated by the set of circle inversions $s_L$ with $L \perp \hat g$.
	\end{enumerate}
\end{proposition}